\section{Infinite-volume gaps and ground states of primitive sectors}

\begin{definition}[Insertion expectation]
    \label{def:primitive_insertion_expectation}
    \lean{MPSTensor.observableInsertionExpectation}
    \leanok
    For a tensor $A$, a matrix $\rho$ with nonzero trace, and an observable
    $X$ on $k$ sites, write
    $E_X(Y)=\sum_{\sigma,\tau}X_{\sigma\tau}A^\tau Y(A^\sigma)^*$ and
    $\eta_A(X)=\operatorname{tr}(E_X(\rho))/\operatorname{tr}\rho$.
    This is the local expectation appearing in
    \cite[Section~6]{Nachtergaele1996Spectral}.
\end{definition}

\begin{definition}[Observable in an interior interval]
    \label{def:primitive_bulk_observable}
    \lean{MPSTensor.bulkObservable}
    \leanok
    For an observable $X$ on $k$ sites and nonnegative integers $\ell,r$,
    let $X_{\ell,r}=I_\ell\otimes X\otimes I_r$ on $N=\ell+k+r$ sites,
    with the tensor factors ordered from left to right.
\end{definition}

\begin{lemma}[Boundary compression of an insertion]
    \label{thm:primitive_insertion_boundary_identity}
    \lean{MPSTensor.adjoint_groundSpaceMapES_observable_groundSpaceMapES,
        MPSTensor.physicalObservableTransfer_bulkObservable,
        MPSTensor.adjoint_groundSpaceMapES_bulkObservable_groundSpaceMapES}
    \leanok
    Let $\Gamma_n$ be the Hilbert-space boundary map and let
    $\mathcal R$ denote Gram reshuffling. Then
    \begin{align}
        \Gamma_k^*X\Gamma_k          & =\mathcal R(E_X),           &
        E_{X_{\ell,r}}               & =E^\ell E_XE^r,               \\
        \Gamma_N^*X_{\ell,r}\Gamma_N & =\mathcal R(E^\ell E_XE^r).
    \end{align}
\end{lemma}
\begin{proof}\leanok
    \uses{def:ground_space_map_es, def:gram_reshuffle, def:primitive_bulk_observable}
    Pair both sides of the first identity with matrix-unit boundary vectors.
    The resulting finite sums agree term by term.
    Insertion transfers multiply under consecutive tensor products,
    while the identity insertion on $n$ sites is $E^n$.
    These facts give the remaining identities.
\end{proof}

\begin{lemma}[Primitive insertion limit]
    \label{thm:primitive_insertion_transfer_limit}
    \lean{MPSTensor.fixedPointProj_physicalObservableTransfer_fixedPointProj,
        MPSTensor.IsPrimitiveMPS.transferMap_pow_tendsto_fixedPointProj,
        MPSTensor.IsPrimitiveMPS.physicalObservableTransfer_sandwich_tendsto}
    \leanok
    Let $A$ be a normalized primitive tensor with invariant matrix $\rho$,
    and let $P_\rho(Y)=\operatorname{tr}(Y)\rho/\operatorname{tr}\rho$.
    For every local observable $X$,
    $P_\rho E_XP_\rho=\eta_A(X)P_\rho$ and $E^n\to P_\rho$ in operator norm.
    Along any filter on which $\ell,r\to\infty$,
    $E^\ell E_XE^r\to\eta_A(X)P_\rho$.
\end{lemma}
\begin{proof}\leanok
    \uses{def:primitive_insertion_expectation, thm:primitive_fnw_remainder_power}
    Substitution into the rank-one projection gives the first identity.
    For $n\geq1$, the decomposition
    $E^n=P_\rho+(E-P_\rho)^n$ and decay of the complementary powers
    give the second. Continuity of multiplication gives the insertion limit.
\end{proof}

\begin{lemma}[Limiting boundary compression]
    \label{thm:primitive_insertion_gram_limit}
    \lean{MPSTensor.IsPrimitiveMPS.groundSpaceMapES_bulkObservable_compression_tendsto}
    \leanok
    Under the preceding hypotheses, with $K_\infty=\mathcal R(P_\rho)$,
    $\Gamma_N^*X_{\ell,r}\Gamma_N\to\eta_A(X)K_\infty$ as both
    free intervals tend to infinity.
\end{lemma}
\begin{proof}\leanok
    \uses{thm:primitive_insertion_boundary_identity, thm:primitive_insertion_transfer_limit}
    Apply the continuous linear reshuffling map to the insertion limit.
\end{proof}

\begin{lemma}[Scalar compression on the entire primitive ground space]
    \label{thm:primitive_bulk_ground_compression}
    \lean{MPSTensor.IsPrimitiveMPS.bulkObservable_groundSpace_compression_tendsto_scalar,
        MPSTensor.IsPrimitiveMPS.bulkObservable_groundSpace_compression_tendsto_zero,
        MPSTensor.IsPrimitiveMPS.bulkObservable_groundSpace_projection_eventually_near_expectation}
    \leanok
    Suppose in addition that $\rho$ is positive definite. Let $G_N$ be the
    open-boundary MPS space and $P_N$ its orthogonal projection. Then
    \begin{align}
        \|P_NX_{\ell,r}P_N-\eta_A(X)P_N\| & \longrightarrow0.
    \end{align}
    In particular, centered observables satisfy
    $\|P_NX_{\ell,r}P_N\|\to0$. For every $\varepsilon>0$, eventually
    $\|P_NX_{\ell,r}\psi-\eta_A(X)\psi\|<\varepsilon$ simultaneously
    for all $\psi\in G_N$ with $\|\psi\|\leq1$.
\end{lemma}
\begin{proof}\leanok
    \uses{thm:primitive_insertion_gram_limit, thm:primitive_ground_space_gram_tendsto_fixed_point,
        thm:primitive_ground_space_map_eventual_inverse_gram_bound,
        thm:injective_boundary_compression_norm}
    Subtract $\eta_A(X)\Gamma_N^*\Gamma_N$ from the limiting boundary
    compression. The difference tends to zero. The boundary maps are
    eventually injective and their inverse Gram operators remain bounded.
    The compression estimate therefore gives the operator-norm limit.
    Apply the resulting operator to vectors of norm at most one.
\end{proof}

\begin{lemma}[Centered projection decay and local expectation convergence]
    \label{thm:primitive_bulk_ground_vector_limits}
    \lean{MPSTensor.IsPrimitiveMPS.bulkObservable_groundSpace_projection_norm_sq_tendsto_zero,
        MPSTensor.IsPrimitiveMPS.bulkObservable_groundState_expectation_tendsto}
    \leanok
    With the preceding hypotheses, let $\psi_N$ be any varying family of
    vectors eventually belonging to $G_N$. If $\|\psi_N\|\leq1$ and
    $\eta_A(X)=0$, then $\|P_NX_{\ell,r}\psi_N\|^2\to0$.
    If instead $\|\psi_N\|=1$, then
    $\langle\psi_N,X_{\ell,r}\psi_N\rangle\to\eta_A(X)$.
    These assertions concern one primitive sector; the general pure-GVBS
    infinite-volume statement is discussed in
    \cite{gap:nachtergaele96_infinite_volume_ground_projection}.
\end{lemma}
\begin{proof}\leanok
    \uses{thm:primitive_bulk_ground_compression, thm:injective_boundary_compression_norm}
    The projection norm is bounded by the centered compression norm.
    Squaring preserves its zero limit. For unit vectors, the scalar
    expectation deviation is bounded by the scalar compression norm.
\end{proof}

\begin{lemma}[Positive observable insertions]
    \label{thm:primitive_insertion_positive}
    \lean{MPSTensor.physicalObservableTransfer_posSemidef}
    \leanok
    For any tensor $A$, positive observable $X$ on $k$ sites, and positive
    virtual matrix $\rho$, the insertion $E_X(\rho)$ is positive.
    This is the positivity of the local GVBS contraction in
    \cite[Section~3, equations (3.1)--(3.2b)]{Nachtergaele1996Spectral}.
\end{lemma}
\begin{proof}\leanok
    \uses{def:primitive_insertion_expectation}
    Let $B_{a,(\sigma,b)}=(A^\sigma)_{ab}$, where $\sigma$ runs over
    words of length $k$. The insertion is
    $E_X(\rho)=B(X^{\mathsf T}\otimes\rho)B^*$.
    Transposition, tensor products, and congruence preserve positivity.
\end{proof}

\begin{lemma}[Local state properties and consistency]
    \label{thm:primitive_insertion_local_state}
    \lean{MPSTensor.observableInsertionExpectation_one,
        MPSTensor.observableInsertionExpectation_nonneg,
        MPSTensor.observableInsertionExpectation_add,
        MPSTensor.observableInsertionExpectation_smul,
        MPSTensor.observableInsertionExpectation_appendObservable_one,
        MPSTensor.observableInsertionExpectation_one_appendObservable,
        MPSTensor.IsPrimitiveMPS.observableInsertionExpectation_one,
        MPSTensor.IsPrimitiveMPS.observableInsertionExpectation_nonneg}
    \leanok
    Fix $A$ and a virtual matrix $\rho$ with nonzero trace. For each
    interval length $k$, the functional $\eta_A$ is complex linear.
    If $\rho\geq0$, then $\eta_A(X)\geq0$ for $X\geq0$.
    If $E(\rho)=\rho$, then $\eta_A(I_k)=1$ and adjoining identity
    factors on the right preserves expectations. If
    $\sum_i(A^i)^*A^i=I$, adjoining identity factors on the left
    preserves expectations. Thus a normalized primitive tensor with
    its positive invariant matrix defines compatible states on finite
    intervals, as in
    \cite[Section~3, equations (3.1)--(3.2b)]{Nachtergaele1996Spectral}.
\end{lemma}
\begin{proof}\leanok
    \uses{def:primitive_insertion_expectation, thm:primitive_insertion_positive}
    Linearity follows from linearity of the insertion and the trace.
    The positive numerator and denominator give positivity.
    At a fixed point, $E^m(\rho)=\rho$, so the identity insertion
    $E_{I_m}=E^m$ gives normalization. Consecutive insertions compose,
    giving $E_{X\otimes I_m}=E_XE^m$ and
    $E_{I_m\otimes X}=E^mE_X$. Invariance removes the right factor,
    while trace preservation removes the left factor after taking the
    trace. The primitive assumptions include positivity, invariance,
    nonzero trace, and trace preservation.
\end{proof}

\begin{lemma}[Limiting norm of a primitive-sector excitation]
    \label{thm:primitive_bulk_excitation_norm_limit}
    \lean{MPSTensor.IsPrimitiveMPS.bulkObservable_groundState_norm_sq_tendsto}
    \leanok
    Let $A$ be a normalized primitive tensor with positive-definite invariant
    matrix $\rho$. Let $X$ act on $k$ sites and let $\ell,r\to\infty$
    along any filter. For any family of unit vectors $\psi_N$ eventually
    belonging to the open-boundary MPS space on $N=\ell+k+r$ sites,
    $\|X_{\ell,r}\psi_N\|^2\to\operatorname{Re}\eta_A(X^*X)$.
\end{lemma}
\begin{proof}\leanok
    \uses{thm:primitive_bulk_ground_vector_limits, def:primitive_bulk_observable}
    The interior inclusion preserves products and adjoints, so
    $\langle\psi_N,(X^*X)_{\ell,r}\psi_N\rangle
        =\|X_{\ell,r}\psi_N\|^2$.
    Apply local expectation convergence to $X^*X$ and take real parts.
\end{proof}

\begin{lemma}[Primitive-sector commutator gap reduction]
    \label{thm:primitive_bulk_commutator_gap_reduction}
    \lean{MPSTensor.IsPrimitiveMPS.commutator_gap_of_bulkObservable_energy_limit}
    \leanok
    Let $A$ be a normalized primitive tensor with positive-definite invariant
    matrix $\rho$, and let $X$ be a fixed observable with $\eta_A(X)=0$.
    Consider a sequence $N_n=\ell_n+k+r_n$ with both $\ell_n,r_n\to\infty$.
    Suppose the operators $H_n$ are eventually positive, have kernel equal
    to the open-boundary MPS space, and have a common norm gap
    $\gamma\geq0$ on its orthogonal complement.
    Let $\psi_n$ be unit vectors eventually in that ground space.
    If
    \begin{align}
        \operatorname{Re}\langle X_{\ell_n,r_n}^*
        [H_n,X_{\ell_n,r_n}]\psi_n,\psi_n\rangle & \longrightarrow e,
    \end{align}
    then $\gamma\operatorname{Re}\eta_A(X^*X)\leq e$.
    The energy limit is an explicit hypothesis here; the passage to the
    infinite-volume state in \cite[Section~6]{Nachtergaele1996Spectral}
    remains separate, as discussed in
    \cite{gap:nachtergaele96_infinite_volume_ground_projection}.
\end{lemma}
\begin{proof}\leanok
    \uses{thm:primitive_bulk_excitation_norm_limit, thm:primitive_bulk_ground_vector_limits,
        thm:commutator_gap_projection_limit}
    The excitation norm squared tends to $\operatorname{Re}\eta_A(X^*X)$,
    while the squared norm of its full ground-space projection tends to zero.
    Take limits in the finite-volume commutator gap inequality.
\end{proof}

\begin{lemma}[Interval inclusions and interior observable placement]
    \label{thm:local_interval_bulk_observable}
    \lean{SpinChain.intervalCoordinates_localInclusion_eq_bulkObservable}
    \leanok
    With $I,J$ as above,
    $C_{a-b,b+N+c}(\iota_{I,J}(X))=I_b\otimes C_{a,N}(X)\otimes I_c$.
    Thus the integer-interval inclusion agrees with the interior placement
    used in the local GVBS expectations of
    \cite[Section~3, equations (3.1)--(3.2b)]{Nachtergaele1996Spectral}.
\end{lemma}
\begin{proof}\leanok
    \uses{thm:local_interval_coordinate_inclusion, def:primitive_bulk_observable}
    Write $\sigma=(\sigma_1,\sigma_2,\sigma_3)$ and $\tau=(\tau_1,\tau_2,\tau_3)$
    for the restrictions of configurations to the first $b$ sites, the middle
    $N$ sites, and the last $c$ sites. The interval inclusion formula gives
    \begin{align}
        C_{a-b,b+N+c}(\iota_{I,J}(X))_{\sigma,\tau}
        =\delta_{\sigma_1,\tau_1}\,C_{a,N}(X)_{\sigma_2,\tau_2}\,
        \delta_{\sigma_3,\tau_3},
        \notag
    \end{align}
    which is the entry of $I_b\otimes C_{a,N}(X)\otimes I_c$.
\end{proof}

\begin{lemma}[Norm of a finite local state]
    \label{thm:primitive_insertion_state_norm}
    \lean{MPSTensor.observableInsertionExpectationₗ,
        MPSTensor.observableInsertionExpectationPositive,
        MPSTensor.observableInsertionExpectationPositive_apply,
        MPSTensor.norm_observableInsertionExpectationPositive,
        MPSTensor.norm_observableInsertionExpectation_le,
        MPSTensor.observableInsertionExpectation_conjTranspose}
    \leanok
    Let $\rho\geq0$ satisfy $E_A(\rho)=\rho$ and
    $\operatorname{tr}\rho\neq0$. For each finite interval, the insertion
    expectation $\eta_A$ is a positive continuous linear functional of
    operator norm one. Consequently
    \begin{align}
        |\eta_A(X)| & \leq\|X\|,
                    & \eta_A(X^\dagger) & =\overline{\eta_A(X)}.
    \end{align}
    The adjoint identity requires only $\rho\geq0$; the norm-one assertion
    also uses invariance and nonzero trace. Primitivity is unnecessary.
    This is the finite local state of
    \cite[Section~3, equations (3.1)--(3.2b)]{Nachtergaele1996Spectral}.
\end{lemma}
\begin{proof}\leanok
    \uses{thm:primitive_insertion_local_state, thm:primitive_insertion_positive}
    Positivity and complex linearity give a positive functional on the
    finite matrix algebra. Its norm is its value at the identity, which is
    one. A positive complex-linear functional preserves adjoints.
\end{proof}

\begin{definition}[Finite open interaction sum in matrix coordinates]
    \label{def:finite_open_interaction_matrix}
    \lean{MPSTensor.openInteractionMatrix,
        MPSTensor.openInteractionHamiltonianES_eq_toEuclideanLin_openInteractionMatrix,
        MPSTensor.toEuclideanCLM_openInteractionMatrix_parentInteraction}
    \leanok
    For an arbitrary range-$R$ matrix $h$, let
    $H_N(h)=\sum_{a=0}^{N-R}\iota_a(h)$, where $\iota_a$ places $h$
    on the sites $a,\ldots,a+R-1$ of an $N$-site chain, with the sum
    empty when $N<R$. For $R\geq1$ and $N\geq R$, its Euclidean matrix
    action is the open interaction Hamiltonian. If $h$ is the complement
    projection of $G_R(A)$, it is the canonical open parent Hamiltonian.
    This is the nonwrapping interaction sum in
    \cite[equation (3.12)]{Nachtergaele1996Spectral}.
\end{definition}
\begin{proof}\leanok
    \uses{def:open_parent_hamiltonian_es}
    The admissible starts are exactly $0,\ldots,N-R$. Each matrix inclusion
    acts as the corresponding Euclidean local interaction, so their sums
    agree. Substitute the local complement projection for the canonical case.
\end{proof}

\begin{lemma}[Commutation on disjoint intervals]
    \label{thm:disjoint_interval_observables_commute}
    \lean{MPSTensor.chainWindowOperator_commute_of_disjoint}
    \leanok
    Observables supported on two disjoint nonwrapping intervals of a finite
    chain commute. No self-adjointness assumption is required.
\end{lemma}
\begin{proof}\leanok
    Operators on the two tensor factors commute; each acts as the identity
    on the support of the other.
\end{proof}

\begin{definition}[The fixed local commutator observable]
    \label{def:fixed_local_commutator_observable}
    \lean{MPSTensor.localCommutatorObservable}
    \leanok
    Let $R,k\geq1$, put $t=R-1$ and $w=2t+k$, and let $X$ be supported
    on $k$ sites. On $w$ sites set $B=I_t\otimes X\otimes I_t$ and
    \begin{align}
        Y_h(X) & =B^\dagger\bigl(H_w(h)B-BH_w(h)\bigr).
    \end{align}
    Neither positivity nor self-adjointness of $h$ is needed for this
    definition. It is the fixed local expression used in the limiting
    argument of \cite[Section~6]{Nachtergaele1996Spectral}.
\end{definition}

\begin{lemma}[Finite-range commutator localization]
    \label{thm:finite_open_commutator_locality}
    \lean{MPSTensor.openInteractionMatrix_commutator_chainWindowOperator,
        MPSTensor.openInteractionMatrix_commutator_bulkObservable,
        MPSTensor.chainWindowOperator_adjoint_mul_openInteractionMatrix_commutator,
        MPSTensor.bulkObservable_adjoint_mul_openInteractionMatrix_commutator}
    \leanok
    Let $R,k\geq1$ and $X_N=I_\ell\otimes X\otimes I_r$, where
    $\ell,r\geq t=R-1$ and $N=\ell+k+r$. Then
    \begin{align}
        [H_N(h),X_N]
         & =I_{\ell-t}\otimes[H_{2t+k}(h),B]\otimes I_{r-t}, \\
        X_N^\dagger[H_N(h),X_N]
         & =I_{\ell-t}\otimes Y_h(X)\otimes I_{r-t}.
    \end{align}
    Equivalently, both identities hold for an arbitrary placement of the
    enlarged interval inside a longer chain. These are locality identities;
    no spectral-gap assumption is made.
\end{lemma}
\begin{proof}\leanok
    \uses{def:finite_open_interaction_matrix,
        def:fixed_local_commutator_observable,
        thm:disjoint_interval_observables_commute, thm:bulk_observable_algebra}
    Terms disjoint from the support of $X$ have zero commutator. The remaining
    terms are precisely the translates in the enlarged interval. Reindex
    their sum and use multiplicativity of the interval inclusion. Multiplying
    by $X_N^\dagger$ gives the second identity.
\end{proof}

\begin{lemma}[Primitive-sector local commutator energy limit]
    \label{thm:primitive_local_commutator_energy_limit}
    \lean{MPSTensor.IsPrimitiveMPS.localCommutatorObservable_expectation_tendsto}
    \leanok
    Let $A$ be a normalized primitive tensor with positive-definite invariant
    matrix $\rho$, let $h$ act on $R>0$ sites, and let $X$ act on $k>0$
    sites. Put $t=R-1$, $w=2t+k$, $B=I_t\otimes X\otimes I_t$, and
    $Y=B^*[H_w(h),B]$. On chains of length $N_n=2n+w$, place $X$ at
    start $n+t$. For any eventually unit vectors $\psi_n$ in the
    open-boundary MPS space,
    $\operatorname{Re}\langle X_n^*[H_{N_n}(h),X_n]\psi_n,\psi_n\rangle
        \to\operatorname{Re}\eta_A(Y)$.
\end{lemma}
\begin{proof}\leanok
    \uses{thm:primitive_bulk_ground_vector_limits, thm:finite_open_commutator_locality}
    The commutator energy operator is the inclusion of $Y$ starting at $n$.
    Apply the primitive local expectation limit to $Y$.
    Conjugate symmetry of the inner product preserves its real part.
\end{proof}

\begin{lemma}[Eventually normalized primitive boundary vectors]
    \label{thm:primitive_unit_boundary_approximants}
    \lean{MPSTensor.IsPrimitiveMPS.exists_eventually_unit_groundSpaceES}
    \leanok
    Let $A$ and $\rho$ have the preceding properties. Along any filter with
    $N_i\to\infty$, there exist vectors $\psi_i$ in the open-boundary MPS
    space at every length $N_i$, whose norms are eventually one.
\end{lemma}
\begin{proof}\leanok
    \uses{thm:primitive_ground_space_map_eventual_inverse_gram_bound,
        def:ground_space_map_es}
    Choose a nonzero matrix-unit boundary vector $u$ and set
    $\psi_i=\|\Gamma_{N_i}u\|^{-1}\Gamma_{N_i}u$, with zero inverse
    interpreted as zero. The boundary map is eventually injective, so these
    vectors are eventually normalized; their ground-space membership holds
    at every length.
\end{proof}

\begin{lemma}[Primitive-sector local commutator inequality]
    \label{thm:primitive_local_commutator_gap}
    \lean{MPSTensor.IsPrimitiveMPS.localCommutator_gap_of_open_gap}
    \leanok
    Let $A$ be a normalized primitive tensor with positive-definite invariant
    matrix $\rho$, and let $X$ act on $k>0$ sites with $\eta_A(X)=0$.
    Fix $R>0$ and let $h=I-P_{G_R(A)}$. Suppose that, on the expanding
    chains $N_n=2n+2(R-1)+k$, the canonical open Hamiltonians eventually
    have kernel $G_{N_n}(A)$ and a common norm gap $\gamma\geq0$.
    Suppose also that a family of unit vectors eventually belongs to these
    ground spaces. With $B=I_{R-1}\otimes X\otimes I_{R-1}$,
    \begin{align}
        \gamma\operatorname{Re}\eta_A(X^*X)
         & \leq \operatorname{Re}\eta_A\bigl(B^*[H_{2(R-1)+k}(h),B]\bigr).
    \end{align}
    This is a statement about finite-interval primitive-sector expectations.
    Their extension to the quasi-local algebra and the GNS Hamiltonian remains
    separate, as discussed in
    \cite{gap:nachtergaele96_infinite_volume_ground_projection}.
\end{lemma}
\begin{proof}\leanok
    \uses{thm:primitive_local_commutator_energy_limit,
        thm:primitive_bulk_excitation_norm_limit,
        thm:primitive_bulk_ground_vector_limits, thm:commutator_gap_projection_limit,
        thm:finite_open_commutator_locality}
    The excitation norm squared tends to $\operatorname{Re}\eta_A(X^*X)$;
    its full ground-space projection tends to zero. The canonical matrix
    interaction sum is the open parent Hamiltonian. Apply the commutator
    energy limit and take limits in the finite-volume gap inequality.
\end{proof}

\begin{lemma}[A positive local commutator bound at the Wielandt range]
    \label{thm:primitive_local_commutator_wielandt}
    \lean{MPSTensor.IsPrimitiveMPS.exists_pos_localCommutator_gap_of_wielandt_range}
    \leanok
    Let $A$ be a normalized primitive tensor of bond dimension $D>0$ with
    positive-definite invariant matrix $\rho$. At each fixed range
    $R\geq D^4+1$, there exists $\gamma>0$ such that the preceding local
    commutator inequality holds for every observable on every nonempty
    finite interval with $\eta_A(X)=0$.
    The interaction is $h=I-P_{G_R(A)}$.
\end{lemma}
\begin{proof}\leanok
    \uses{thm:normal_from_primitive_posdef, thm:normal_open_gap_every_sufficient_range,
        thm:open_parent_hamiltonian_kernel_ground_space,
        thm:primitive_unit_boundary_approximants, thm:primitive_local_commutator_gap}
    Primitivity and positive definiteness imply normality.
    The dimension bound gives injectivity at $D^4$ and hence at $R-1$.
    The open-chain kernel is therefore the MPS ground space, and the normal
    open-chain gap theorem gives one positive constant independent of
    length. Choose eventually normalized boundary vectors and apply the
    preceding inequality.
\end{proof}

\begin{lemma}[Equality of open parent-interaction kernels]
    \label{thm:open_parent_interaction_kernel_equal}
    \lean{MPSTensor.IsParentInteraction.%
        ker_openInteractionHamiltonianES_eq_ker_openParentHamiltonianES}
    \leanok
    Let $h\geq0$ act on $R>0$ sites with kernel $G_R(A)$, and let
    $q=I-P_{G_R(A)}$. For every chain length $N$,
    $\ker H_N(h)=\ker H_N(q)$.
\end{lemma}
\begin{proof}\leanok
    \uses{thm:open_positive_parent_gap}
    There are $\kappa,C>0$ with
    $\kappa H_N(q)\leq H_N(h)\leq C H_N(q)$ at every length.
    For positive operators, either bound gives the corresponding reverse
    inclusion of kernels. Their combination gives equality.
\end{proof}

\begin{lemma}[Primitive-sector commutator bound for an open interaction]
    \label{thm:primitive_local_open_interaction_commutator_gap}
    \lean{MPSTensor.IsPrimitiveMPS.localCommutator_gap_of_openInteraction_gap}
    \leanok
    Let $A$ be a normalized primitive tensor with positive-definite invariant
    matrix $\rho$, let $h\geq0$ act on $R>0$ sites, and let $X$ act on
    $k>0$ sites with $\eta_A(X)=0$. Suppose the open interaction sums
    on chains $N_n=2n+2(R-1)+k$ eventually have kernel $G_{N_n}(A)$
    and common norm gap $\gamma\geq0$. Suppose that a family of unit
    vectors eventually belongs to these ground spaces.
    With $B=I_{R-1}\otimes X\otimes I_{R-1}$,
    \begin{align}
        \gamma\operatorname{Re}\eta_A(X^*X)
         & \leq\operatorname{Re}\eta_A\bigl(B^*[H_{2(R-1)+k}(h),B]\bigr).
    \end{align}
\end{lemma}
\begin{proof}\leanok
    \uses{thm:primitive_bulk_commutator_gap_reduction,
        thm:primitive_local_commutator_energy_limit, thm:finite_open_commutator_locality}
    The open matrix sum represents the open interaction Hamiltonian.
    Its commutator energy tends to the expectation of the fixed local patch.
    Apply the primitive boundary commutator reduction, which derives the
    excitation norm and full ground-space projection limits from the tensor.
\end{proof}

\begin{lemma}[Positive parent interactions satisfy a primitive-sector local gap]
    \label{thm:primitive_local_positive_parent_commutator_gap}
    \lean{MPSTensor.IsPrimitiveMPS.exists_pos_localCommutator_gap_of_isParentInteraction}
    \leanok
    Let $A$ be a normalized primitive tensor of bond dimension $D>0$ with
    positive-definite invariant matrix $\rho$. Let $h\geq0$ act on
    $R\geq D^4+1$ sites with kernel $G_R(A)$. There exists $\gamma>0$
    such that the preceding local commutator bound holds for every
    observable on every nonempty finite interval with $\eta_A(X)=0$.
    This concerns one primitive sector and its finite-interval expectations;
    the quasi-local, GNS, and pure multiblock conclusions remain separate
    \cite{gap:nachtergaele96_infinite_volume_ground_projection}.
\end{lemma}
\begin{proof}\leanok
    \uses{thm:normal_from_primitive_posdef, thm:normal_open_gap_every_sufficient_range,
        thm:open_positive_parent_gap, thm:open_parent_interaction_kernel_equal,
        thm:open_parent_hamiltonian_kernel_ground_space,
        thm:primitive_unit_boundary_approximants,
        thm:primitive_local_open_interaction_commutator_gap}
    The canonical open Hamiltonian has a uniform positive gap and the MPS
    ground-space kernel. Positive local comparison transfers the gap and
    preserves that kernel. Choose eventually normalized boundary vectors
    and apply the preceding local commutator inequality.
\end{proof}

\begin{lemma}[Compression between two injective ranges]
    \label{thm:parent_cross_range_compression}
    \lean{ContinuousLinearMap.norm_starProjection_comp_starProjection_le_inverseGram,
        ContinuousLinearMap.tendsto_starProjection_comp_starProjection_zero_of_inverseGram_bounds}
    \leanok
    Let $T\colon E\to F$ and $S\colon G\to F$ be injective bounded linear
    maps between real or complex Hilbert spaces, with $E$ and $G$
    finite-dimensional. Write $K_T=(T^\dagger T)^{-1}$,
    $K_S=(S^\dagger S)^{-1}$, and let $P_T,P_S$ be the orthogonal
    projections onto their ranges. For every bounded operator $X$ on $F$,
    \begin{align}
        \lVert P_T X P_S\rVert
         & \leq\lVert K_T\rVert^{1/2}
        \lVert T^\dagger X S\rVert\lVert K_S\rVert^{1/2}.
    \end{align}
    Consequently, along any filter, if these maps are eventually injective,
    their inverse Gram norms are eventually bounded by constants, and
    $\lVert T^\dagger X S\rVert\to0$, then
    $\lVert P_T X P_S\rVert\to0$. The Hilbert spaces and maps may vary
    with the index.
\end{lemma}

\begin{proof}
    \leanok
    Since $P_T=TK_TT^\dagger$ and $P_S=SK_SS^\dagger$,
    factor the compression as
    $(TK_T)(T^\dagger X S)(K_SS^\dagger)$.
    The identities $\lVert TK_T\rVert=\lVert K_T\rVert^{1/2}$ and
    $\lVert K_SS^\dagger\rVert=\lVert K_S\rVert^{1/2}$ give the bound.
    The convergence assertion follows by comparison with
    $\sqrt C\lVert T^\dagger X S\rVert\sqrt D$, where $C,D$ are
    eventual bounds for the two inverse Gram norms.
\end{proof}

\begin{definition}[MPS expectations on finite regions]
    \label{def:mps_finite_region_expectation}
    \lean{MPSTensor.localObservableExpectation}
    \leanok
    For a finite region $\Lambda\subset\mathbb Z$, include an observable
    $X\in\mathcal A_\Lambda$ into a consecutive interval $I(a,N)$
    containing $\Lambda$ and put
    \begin{align}
        \eta_{A,\Lambda}(X)
         & =\eta_A\bigl(C_{a,N}(\iota_{\Lambda,I(a,N)}X)\bigr).
    \end{align}
    To define the functional before proving independence, one may use
    $I(-r,2r+1)$ with $r=\max_{x\in\Lambda}|x|$ and $r=0$ when
    $\Lambda=\varnothing$. This gives a complex-linear functional without
    a normalization assumption on $A$ or $\rho$.
\end{definition}

\begin{lemma}[Independence and compatibility of finite-region expectations]
    \label{thm:mps_finite_region_expectation_compatible}
    \lean{MPSTensor.localObservableExpectation_eq_of_subset_interval,
        MPSTensor.localObservableExpectation_localInclusion}
    \leanok
    Suppose $\sum_i A_i^\dagger A_i=I$ and $E_A(\rho)=\rho$.
    Then the preceding definition is independent of the containing interval.
    For $\Lambda\subseteq\Gamma$ and $X\in\mathcal A_\Lambda$,
    $\eta_{A,\Gamma}(\iota_{\Lambda,\Gamma}X)=\eta_{A,\Lambda}(X)$.
    Neither positivity nor nonzero trace of $\rho$ is required for these
    algebraic compatibility identities.
\end{lemma}
\begin{proof}\leanok
    \uses{def:mps_finite_region_expectation,thm:local_interval_bulk_observable,thm:primitive_insertion_local_state,lem:qca_local_inclusion_trans}
    Identity extension on the left preserves expectation by trace preservation;
    identity extension on the right preserves expectation by $E_A(\rho)=\rho$.
    Include two containing intervals in a common larger interval to compare
    their values. Composition of local inclusions gives finite-region compatibility.
\end{proof}

\begin{definition}[The compatible local MPS state]
    \label{def:mps_compatible_local_state}
    \lean{MPSTensor.localMPSState, MPSTensor.localMPSState_functional_interval}
    \leanok
    Suppose $\sum_i A_i^\dagger A_i=I$, $\rho\geq0$,
    $E_A(\rho)=\rho$, and $\operatorname{tr}\rho\neq0$.
    The functionals $\eta_{A,\Lambda}$ form a compatible family of
    normalized positive contractive local states. In particular,
    $\eta_{A,I(a,N)}(X)=\eta_A(C_{a,N}X)$ for every $N\geq0$.
    This is the compatible local interpretation of
    \cite[Section~3, equations (3.1)--(3.2b)]{Nachtergaele1996Spectral}.
\end{definition}
\begin{proof}\leanok
    \uses{thm:mps_finite_region_expectation_compatible,thm:primitive_insertion_state_norm,def:qca_compatible_local_state,thm:primitive_insertion_local_state}
    The inclusion and coordinate maps preserve identities, products, and
    adjoints. Positivity and normalization therefore follow from their
    finite-interval counterparts. Star-algebra homomorphisms between
    complex C$^*$-algebras are contractive, so
    $|\eta_{A,\Lambda}(X)|\leq\|X\|$.
\end{proof}

\begin{definition}[The completed quasi-local MPS expectation]
    \label{def:mps_quasi_local_expectation}
    \lean{MPSTensor.quasiLocalExpectation}
    \leanok
    Under the preceding hypotheses and $d>0$, let $\omega_A$ be the
    continuous linear extension of the compatible local MPS state to the
    completed quasi-local observable algebra.
\end{definition}

\begin{lemma}[Local evaluation and state properties of the MPS extension]
    \label{thm:mps_quasi_local_expectation_state}
    \lean{MPSTensor.quasiLocalExpectation_quasiLocalObservable,
        MPSTensor.quasiLocalExpectation_interval, MPSTensor.quasiLocalExpectation_isState}
    \leanok
    Under the hypotheses of the preceding definition, the restriction of
    $\omega_A$ to $\mathcal A_\Lambda$ is $\eta_{A,\Lambda}$.
    Thus $\omega_A(X)=\eta_A(C_{a,N}X)$ for an observable on $I(a,N)$,
    including $N=0$. The extension satisfies
    $\|\omega_A\|=1$, $\omega_A(I)=1$, and $\omega_A(X^*X)\geq0$
    for every quasi-local observable $X$. No primitivity or purity is asserted.
\end{lemma}
\begin{proof}\leanok
    \uses{def:mps_quasi_local_expectation, def:mps_compatible_local_state,
        thm:qca_compatible_state_evaluation, thm:qca_compatible_state_extension}
    Apply the completion theorem for compatible local states. Its finite-region
    evaluation formula gives the asserted restriction, and interval evaluation
    gives the original trace insertion formula.
\end{proof}

\begin{lemma}[Translation preserves interval coordinates]
    \label{thm:local_interval_translation_coordinates}
    \lean{SpinChain.translateRegion_finiteChainRegion,
        SpinChain.intervalCoordinates_localTranslation}
    \leanok
    For $a,b\in\mathbb Z$ and $N\geq0$, translation by $b$ sends
    $I(a,N)$ to $I(a+b,N)$. If $T_b$ denotes the induced local
    star-algebra equivalence, then
    $C_{a+b,N}(T_bX)=C_{a,N}(X)$.
    The equality includes the empty interval.
\end{lemma}
\begin{proof}\leanok
    \uses{def:local_interval_coordinates,def:qca_finite_local_translation,def:qca_translate_region,def:qca_local_inclusion}
    The site $a+i$ is sent to $a+b+i$ and retains its numbered coordinate
    $i$. Hence the two interval coordinate maps give the same configurations
    and matrix entries. The local inclusion identifying the translated region
    with $I(a+b,N)$ has empty complement.
\end{proof}

\begin{lemma}[Translation invariance of finite-region MPS expectations]
    \label{thm:mps_finite_region_translation_invariance}
    \lean{MPSTensor.localObservableExpectation_localTranslation}
    \leanok
    Suppose $\sum_i A_i^\dagger A_i=I$ and $E_A(\rho)=\rho$.
    For every finite region $\Lambda\subset\mathbb Z$, displacement
    $b\in\mathbb Z$, and observable $X\in\mathcal A_\Lambda$,
    $\eta_{A,\Lambda+b}(T_bX)=\eta_{A,\Lambda}(X)$.
    Positivity and nonzero trace of $\rho$ are unnecessary for this identity.
\end{lemma}
\begin{proof}\leanok
    \uses{thm:mps_finite_region_expectation_compatible,thm:local_interval_translation_coordinates,lem:qca_finite_local_translation_naturality,def:qca_translate_region}
    Choose an interval containing $\Lambda$. Translation commutes with the
    local inclusion and preserves numbered interval coordinates. Evaluate
    the two local expectations in this interval and its translate.
\end{proof}

\begin{lemma}[Translation invariance of the completed MPS state]
    \label{thm:mps_quasi_local_translation_invariance}
    \lean{MPSTensor.quasiLocalExpectation_quasiLocalTranslation}
    \leanok
    Suppose $d>0$, $\sum_i A_i^\dagger A_i=I$, $\rho\geq0$,
    $E_A(\rho)=\rho$, and $\operatorname{tr}\rho\neq0$.
    For every displacement $b\in\mathbb Z$ and quasi-local observable $X$,
    the completed state satisfies $\omega_A(T_bX)=\omega_A(X)$.
    This proves translation invariance of the local GVBS state in
    \cite[Section~3, equations (3.1)--(3.2b)]{Nachtergaele1996Spectral};
    no purity conclusion is asserted.
\end{lemma}
\begin{proof}\leanok
    \uses{thm:mps_finite_region_translation_invariance,
        thm:mps_quasi_local_expectation_state, def:qca_quasi_local_translation}
    The equality holds on every finite-region algebra. Translation and the
    completed expectation are continuous, so it extends from the dense
    algebraic local algebra to its completion.
\end{proof}

\begin{lemma}[Interval evaluation of the completed MPS state]
    \label{thm:mps_quasi_local_interval_expectation}
    \lean{MPSTensor.quasiLocalExpectation_quasiLocalIntervalObservable}
    \leanok
    Suppose $d>0$, $\sum_i A_i^\dagger A_i=I$, $\rho\geq0$,
    $E_A(\rho)=\rho$, and $\operatorname{tr}\rho\neq0$.
    For every $a\in\mathbb Z$, $k\geq0$, and $X\in M_{d^k}(\mathbb C)$,
    the completed state satisfies $\omega_A(j_{a,k}(X))=\eta_A(X)$.
\end{lemma}
\begin{proof}\leanok
    \uses{thm:mps_quasi_local_expectation_state, thm:quasi_local_interval_observable}
    Apply interval evaluation after mapping the matrix observable to
    the interval algebra. The inverse coordinate maps cancel.
\end{proof}

\begin{lemma}[A primitive-sector quasi-local commutator gap]
    \label{thm:primitive_quasi_local_parent_commutator_gap}
    \lean{MPSTensor.IsPrimitiveMPS.exists_pos_quasiLocalCommutator_gap_of_isParentInteraction}
    \leanok
    Let $d,D>0$, and let $A$ be a normalized primitive tensor with a
    positive-definite invariant matrix $\rho$. Suppose that $h\geq0$ has
    kernel $G_R(A)$ and $R\geq D^4+1$. There is a constant $\gamma>0$
    such that, for every $a\in\mathbb Z$, $k\geq1$, and
    $X\in M_{d^k}(\mathbb C)$ with $\omega_A(j_{a,k}(X))=0$,
    \begin{align}
        \gamma\operatorname{Re}\omega_A\bigl(j_{a,k}(X)^*j_{a,k}(X)\bigr)
        \leq \operatorname{Re}\omega_A\bigl(j_{a-t,w}(Y)\bigr),
        \notag
    \end{align}
    where $t=R-1$, $w=2t+k$, $B=I_t\otimes X\otimes I_t$, and
    $Y=B^*[H_w(h),B]$. The constant is independent of $a$, $k$, and $X$.
    This is the local commutator inequality in the constructed primitive
    quasi-local state, at sufficient range. Identification of all pure
    multiblock GVBS states and construction of the GNS Hamiltonian remain
    separate from this consequence of
    \cite[Theorem~1.2 and Section~6]{Nachtergaele1996Spectral}.
\end{lemma}
\begin{proof}\leanok
    \uses{thm:primitive_local_positive_parent_commutator_gap,
        thm:mps_quasi_local_interval_expectation, thm:quasi_local_interval_observable}
    The interval inclusion preserves products and adjoints. Evaluate both
    sides with the trace insertion formula and apply the finite-interval
    primitive-sector bound.
\end{proof}

\begin{lemma}[Prefix ground projection]
    \label{thm:ground_prefix_projection_identity}
    \lean{MPSTensor.chainWindowOperator_groundSpaceProjection_comp_groundSpaceProjection}
    \leanok
    \uses{def:ground_space_es}
    Let $A$ be an MPS tensor, let $1\leq L\leq N$, and denote its open
    ground projection on $n$ sites by $P_A(n)$. If $Q_A(L,N)$ is $P_A(L)$
    placed on the first $L$ sites and extended by the identity, then
    \begin{align}
        Q_A(L,N)P_A(N)=P_A(N)=P_A(N)Q_A(L,N).
        \notag
    \end{align}
    No injectivity or normalization is assumed.
\end{lemma}
\begin{proof}\leanok
    \uses{thm:open_ground_space_le_hamiltonian_kernel}
    Every open MPS vector satisfies the parent constraint on its prefix.
    Thus the complementary prefix projection annihilates the full ground
    space. Taking adjoints gives the identity in the other order.
\end{proof}

\begin{lemma}[Local cross compression]
    \label{thm:bulk_cross_compression_prefix_bound}
    \lean{MPSTensor.norm_groundSpaceES_bulkObservable_cross_compression_le}
    \leanok
    Let $A,B$ be arbitrary MPS tensors with the same physical dimension.
    Let $X$ act on $k\geq1$ sites, and let $\ell\geq1$ and $r\geq0$.
    Write $N=\ell+k+r$ and let $X_{\ell,r}$ be the placement of $X$ after
    the first $\ell$ sites. Then
    \begin{align}
        \|P_A(N)X_{\ell,r}P_B(N)\|
        \leq \|P_A(\ell)P_B(\ell)\|\,\|X\|.
        \notag
    \end{align}
\end{lemma}
\begin{proof}\leanok
    \uses{thm:ground_prefix_projection_identity,thm:disjoint_interval_observables_commute}
    The observable commutes with both prefix projections. The preceding
    identity therefore gives
    $P_AX_{\ell,r}P_B=P_AQ_AQ_BX_{\ell,r}P_B$.
    The norm of $Q_AQ_B$ is bounded by the norm of the prefix product,
    while the remaining ground projections are contractions.
\end{proof}

\begin{lemma}[Cross-sector decay by prefix locality]
    \label{thm:bulk_cross_compression_decay}
    \lean{MPSTensor.IsPrimitiveMPS.groundSpaceES_starProjection_comp_tendsto_zero_of_inequivalent,
        MPSTensor.IsPrimitiveMPS.bulkObservable_cross_groundSpace_compression_tendsto_zero_of_inequivalent}
    \leanok
    Let $A,B$ be trace-preserving primitive MPS tensors with positive
    definite invariant matrices and positive bond dimensions. Suppose that $A$ and $B$ are inequivalent under a bond
    similarity and a phase, including any identification of equal bond
    dimensions. Then
    \begin{align}
        \|P_A(n)P_B(n)\|\longrightarrow0\qquad(n\longrightarrow\infty).
        \notag
    \end{align}
    For any fixed observable $X$ on $k\geq1$ sites, any filter along which
    $\ell\longrightarrow\infty$, and any choice of $r\geq0$, one also has
    \begin{align}
        \|P_A(\ell+k+r)X_{\ell,r}P_B(\ell+k+r)\|
        \longrightarrow0.
        \notag
    \end{align}
    No growth condition is required of the right interval. This is a
    locality consequence of the overlap estimate in
    \cite[proof of the lemma on distinct local support spaces]{Nachtergaele1996Spectral}.
\end{lemma}
\begin{proof}\leanok
    \uses{thm:inequivalent_block_ground_space_overlap_decay,thm:bulk_cross_compression_prefix_bound}
    The relative overlap estimate bounds the product of the two ground
    projections by an arbitrary positive number at sufficiently large
    length. The second assertion follows by multiplying this bound by
    the fixed norm of the observable and applying the preceding lemma.
\end{proof}

\begin{lemma}[Vanishing error of a finite sum of sector projections]
    \label{thm:sector_projection_sum_error_decay}
    \lean{Submodule.tendsto_norm_sum_starProjection_sub_iSup_zero}
    \leanok
    Let $G_{j,n}$ be a finite family of subspaces of finite-dimensional
    complex Hilbert spaces $E_n$, along an arbitrary filter. Suppose that
    for every $\varepsilon>0$, eventually every pair $i\neq j$ satisfies
    $|\langle x,y\rangle|\leq\varepsilon\|x\|\|y\|$ for
    $x\in G_{i,n}$ and $y\in G_{j,n}$. Write $P_{j,n}$ for their
    orthogonal projections and $P_n$ for the projection onto their joint span.
    Then $\|\sum_jP_{j,n}-P_n\|\to0$.
\end{lemma}
\begin{proof}\leanok
    \uses{thm:ph_sum_block_projectors}
    Let $C_{ij}=1$ for $i\neq j$ and $C_{ii}=0$. The overlap matrix
    $B=\varepsilon C$ satisfies $\|B\|\to0$. The finite-family comparison
    bounds the projection error by $\delta/(1-\delta)$, where
    $\delta=\|B\|/(1-\|B\|)$. This bound tends to zero.
\end{proof}

\begin{lemma}[Combining vanishing sector compressions]
    \label{thm:sector_compression_assembly}
    \lean{Submodule.tendsto_norm_iSup_starProjection_comp_zero_of_sector_compressions}
    \leanok
    Under the preceding asymptotic overlap hypothesis, let $W_n\subseteq E_n$
    have orthogonal projection $Q_n$, and let the operators $X_n$ be
    eventually uniformly bounded. If $\|P_{j,n}X_nQ_n\|\to0$ for every
    sector $j$, then $\|P_nX_nQ_n\|\to0$. The input space $W_n$ may be
    arbitrary and need not be a member of the finite family.
\end{lemma}
\begin{proof}\leanok
    \uses{thm:sector_projection_sum_error_decay}
    Since $\|Q_n\|\leq1$,
    \begin{align}
        \|P_nX_nQ_n\|
         & \leq\|P_n-\textstyle\sum_jP_{j,n}\|\,\|X_n\|
        +\sum_j\|P_{j,n}X_nQ_n\|.
    \end{align}
    Both terms tend to zero, by the projection comparison and finiteness
    of the sector family.
\end{proof}

\begin{lemma}[Centered observables and the full multiblock ground projection]
    \label{thm:multiblock_centered_full_ground_compression}
    \lean{MPSTensor.bulkObservable_iSup_groundSpace_compression_tendsto_zero}
    \leanok
    Let $A_j$ be finitely many normalized primitive tensors with positive
    bond dimensions $D_j$, positive-definite invariant matrices $\rho_j$,
    and pairwise gauge-and-phase inequivalence whenever their bond dimensions
    agree. Fix a sector $\alpha$ and an observable $X$ on $k>0$ sites
    with $\eta_{A_\alpha}(X)=0$. Along an arbitrary filter let
    $\ell_n,r_n\to\infty$ and $N_n=\ell_n+k+r_n$.
    If $P_n$ projects onto $\sum_jG_{N_n}(A_j)$ and $P_{\alpha,n}$
    onto $G_{N_n}(A_\alpha)$, then
    \begin{align}
        \bigl\|P_n(I_{\ell_n}\otimes X\otimes I_{r_n})P_{\alpha,n}\bigr\|
         & \longrightarrow0.
    \end{align}
    This gives the full finite-ground-space projection term in
    \cite[Section~6]{Nachtergaele1996Spectral}. No finite-chain uniqueness
    or purity of a completed state is inferred.
\end{lemma}
\begin{proof}\leanok
    \uses{thm:sector_compression_assembly, thm:primitive_bulk_ground_compression,
        thm:bulk_cross_compression_decay, thm:inequivalent_block_ground_space_overlap_decay,
        thm:bulk_observable_algebra}
    The sector overlaps tend uniformly to zero by inequivalence and
    primitivity. The diagonal compression tends to zero because the
    observable is centered in sector $\alpha$. Every off-diagonal
    compression tends to zero by prefix locality. Interior placement is
    contractive, so the finite-sector combination theorem applies.
\end{proof}

\begin{lemma}[Exact local support of the MPS state]
    \label{thm:mps_state_local_support}
    \lean{MPSTensor.observableInsertionExpectation_eq_zero_of_groundSpaceES_le_ker,
        MPSTensor.quasiLocalExpectation_interval_eq_zero_of_groundSpaceES_le_ker,
        MPSTensor.IsParentInteraction.quasiLocalExpectation_interval_eq_zero}
    \leanok
    Let $A$ be an MPS tensor and let $X$ act on $k$ sites. If
    $G_k(A)\subseteq\ker X$, then $\eta_A(X)=0$ for every virtual
    matrix $\rho$, without normalization or positivity assumptions.
    If $d>0$, $\sum_i A_i^\dagger A_i=I$, $\rho\geq0$,
    $E_A(\rho)=\rho$, and $\operatorname{tr}\rho\neq0$, this gives
    $\omega_A(j_{a,k}(X))=0$ for every $a\in\mathbb Z$.
    In particular, the constructed state has expectation zero for every
    translate of a positive local parent interaction.
    This is the local support property in
    \cite[Theorem~1.1 and Section~3]{Nachtergaele1996Spectral}.
\end{lemma}
\begin{proof}\leanok
    \uses{thm:primitive_insertion_boundary_identity,thm:quasi_local_interval_observable,thm:mps_quasi_local_interval_expectation}
    The hypothesis gives $X\Gamma_k=0$, and hence $\Gamma_k^*X\Gamma_k=0$.
    The latter operator is the Choi reshuffling of the insertion transfer
    map. Its matrix-unit coefficients therefore imply that this map is
    zero. Taking the trace at $\rho$ proves the first assertion; interval
    evaluation proves the assertions for the completed state.
\end{proof}

\begin{lemma}[Vanishing full projection of a sector excitation]
    \label{thm:multiblock_sector_excitation_projection_limit}
    \lean{MPSTensor.bulkObservable_iSup_groundSpace_projection_norm_sq_tendsto_zero}
    \leanok
    Let $(A_j,\rho_j)_{j\in J}$ be a finite family of pairwise inequivalent
    trace-preserving primitive tensors with positive bond dimensions and
    positive-definite invariant matrices. Write
    $G(N)=\sum_jG_{A_j}(N)$ and let $P(N)$ be its orthogonal projection.
    Fix a sector $\alpha$, an observable $X$ on $k\geq1$ sites with
    $\eta_\alpha(X)=0$, and free lengths $\ell,r\longrightarrow\infty$.
    For any varying boundary vectors $\psi\in G_{A_\alpha}(\ell+k+r)$
    with eventually $\|\psi\|\leq1$,
    \begin{align}
        \|P(\ell+k+r)X_{\ell,r}\psi\|^2\longrightarrow0.
        \notag
    \end{align}
    The assertion holds along an arbitrary filter.
\end{lemma}
\begin{proof}\leanok
    \uses{thm:multiblock_centered_full_ground_compression}
    The sector projection fixes $\psi$. Thus the vector norm is bounded
    by the norm of the compression from that sector into the full ground
    space, multiplied by $\|\psi\|$. Apply the compression limit and square.
\end{proof}

\begin{lemma}[Multiblock commutator reduction]
    \label{thm:multiblock_commutator_gap_reduction}
    \lean{MPSTensor.multiblock_commutator_gap_of_bulkObservable_energy_limit}
    \leanok
    Under the preceding tensor hypotheses, suppose that $H_n$ are
    eventually positive operators on chains of length $N_n=\ell_n+k+r_n$,
    with $\ell_n,r_n\longrightarrow\infty$, and eventually
    \begin{align}
        \ker H_n=G(N_n),\qquad
        \gamma\|v\|\leq\|H_nv\|\quad(v\perp\ker H_n),
        \qquad \gamma\geq0.
        \notag
    \end{align}
    Let $\psi_n\in G_{A_\alpha}(N_n)$ be eventually unit vectors, and
    let $B_n=X_{\ell_n,r_n}$ with $\eta_\alpha(X)=0$. If
    \begin{align}
        \operatorname{Re}\langle B_n^*[H_n,B_n]\psi_n,\psi_n\rangle
        \longrightarrow e,
        \notag
    \end{align}
    then
    \begin{align}
        \gamma\operatorname{Re}\eta_\alpha(X^*X)\leq e.
        \notag
    \end{align}
\end{lemma}
\begin{proof}\leanok
    \uses{thm:multiblock_sector_excitation_projection_limit,thm:primitive_bulk_excitation_norm_limit,thm:commutator_gap_projection_limit}
    The distinguished sector lies in the full ground space. Primitive
    convergence gives $\|B_n\psi_n\|^2\to\operatorname{Re}\eta_\alpha(X^*X)$,
    and the preceding theorem removes its full ground component. Apply the
    positive finite-volume gap inequality and pass to the limit.
\end{proof}

\begin{lemma}[Local multiblock commutator inequality]
    \label{thm:multiblock_local_open_interaction_commutator_gap}
    \lean{MPSTensor.multiblock_localCommutator_gap_of_openInteraction_gap}
    \leanok
    Under the same finite-family tensor hypotheses, let $h$ be a positive
    interaction on $R\geq1$ sites. Put $t=R-1$,
    $N_n=(n+t+k)+(t+n)$, and let $H_n$ be its open interaction sum.
    Suppose that these operators eventually have kernel $G(N_n)$ and
    norm gap at least $\gamma\geq0$. For eventually unit
    $\psi_n\in G_{A_\alpha}(N_n)$ and every $k\geq1$ observable $X$
    with $\eta_\alpha(X)=0$,
    \begin{align}
        \gamma\operatorname{Re}\eta_\alpha(X^*X)
        \leq\operatorname{Re}\eta_\alpha(C_h(X)),
        \qquad C_h(X)=B^*[H_{t+k+t},B],\quad B=X_{t,t}.
        \notag
    \end{align}
\end{lemma}
\begin{proof}\leanok
    \uses{thm:multiblock_commutator_gap_reduction,thm:primitive_local_commutator_energy_limit}
    Finite-range locality places the commutator energy on the fixed
    interval of length $t+k+t$. Primitive expectation convergence determines
    its limit, so the preceding reduction applies without a given
    energy limit.
\end{proof}

\begin{lemma}[A common local commutator bound for all primitive sectors]
    \label{thm:multiblock_local_parent_commutator_gap}
    \lean{MPSTensor.exists_pos_multiblock_localCommutator_gap_of_isParentInteraction,
        MPSTensor.exists_pos_multiblock_localCommutator_gap_of_wordTupleSpanTop}
    \leanok
    Let $A_j$ be trace-preserving primitive tensors of positive bond
    dimensions, with positive definite invariant matrices $\rho_j$,
    and let $T=\bigoplus_j\mu_jA_j$, where every $\mu_j\neq0$.
    Assume that the simultaneous word span is the full product of the bond
    matrix algebras at a length $S\geq1$. Let $R\geq S+1$, and let $h$ be
    a positive interaction on $R$ sites with $\ker h=G_T(R)$.
    There exists $\gamma>0$ such that, for every sector $\alpha$ and every
    observable $X$ on any $k\geq1$ sites with $\eta_\alpha(X)=0$,
    \begin{align}
        \gamma\operatorname{Re}\eta_\alpha(X^*X)
        \leq\operatorname{Re}\eta_\alpha(C_h(X)).
        \notag
    \end{align}
    In particular this holds for the canonical parent projection
    $h=I-P_T(R)$. These are finite-interval sector expectation statements;
    they do not assert the construction of a GNS Hamiltonian or purity of
    its state. The range is sufficient and need not be optimal.
    \cite[Theorem~1.2 and Section~6]{Nachtergaele1996Spectral},
    \cite{gap:nachtergaele96_infinite_volume_ground_projection}.
\end{lemma}
\begin{proof}\leanok
    \uses{thm:multiblock_local_open_interaction_commutator_gap,thm:block_open_gap_all_lengths_simultaneous_injectivity,thm:block-word-span-separation,thm:block_open_kernel_simultaneous_injectivity,thm:open_positive_parent_gap,thm:open_parent_interaction_kernel_equal,thm:primitive_unit_boundary_approximants}
    Simultaneous injectivity implies sector inequivalence and identifies
    the open-chain kernel with the sum of the sector spaces at this range.
    The open-chain theorem gives a uniform canonical gap; positivity
    and the prescribed local kernel transfer it to $h$.
    Choose eventual unit boundary vectors independently in each sector
    and apply the preceding local inequality. The same gap constant works
    for every sector and every observable.
\end{proof}

\begin{lemma}[One quasi-local commutator gap in all constructed sectors]
    \label{thm:multiblock_quasi_local_parent_commutator_gap}
    \lean{MPSTensor.exists_pos_multiblock_quasiLocalCommutator_gap_of_isParentInteraction}
    \leanok
    Let $d>0$, and let $A_j$ be finitely many normalized primitive tensors
    with positive bond dimensions and positive-definite invariant matrices
    $\rho_j$. Suppose that the simultaneous word span is full at a length
    $S>0$. Let $T=\bigoplus_j\mu_jA_j$ with $\mu_j\neq0$, and let $h\geq0$
    have kernel $G_R(T)$ at range $R\geq S+1$.
    There is a constant $\gamma>0$ such that, for every sector $\alpha$,
    $a\in\mathbb Z$, $k\geq1$, and $X\in M_{d^k}(\mathbb C)$
    with $\omega_{A_\alpha}(j_{a,k}(X))=0$,
    \begin{align}
        \gamma\operatorname{Re}\omega_{A_\alpha}
        \bigl(j_{a,k}(X)^*j_{a,k}(X)\bigr)
        \leq\operatorname{Re}\omega_{A_\alpha}(j_{a-t,w}(Y)),
        \notag
    \end{align}
    where $t=R-1$, $w=2t+k$, $B=I_t\otimes X\otimes I_t$, and
    $Y=B^*[H_w(h),B]$. The same constant applies to all sectors,
    positions, and centered observables. The states are the completed
    expectations of the given primitive blocks; identification of the full
    pure GVBS ground-state face remains a separate assertion.
\end{lemma}
\begin{proof}\leanok
    \uses{thm:multiblock_local_parent_commutator_gap,
        thm:mps_quasi_local_interval_expectation, thm:quasi_local_interval_observable}
    The common finite-interval bound is expressed in the constructed states
    by interval evaluation. The interval inclusions preserve products and
    adjoints, so the centered hypothesis and both terms agree exactly.
\end{proof}

\begin{definition}[Rectangular observable insertion]
    \label{def:mixed_observable_insertion}
    \lean{MPSTensor.mixedPhysicalObservableTransfer, MPSTensor.mixedPhysicalObservableTransfer_apply, MPSTensor.mixedPhysicalObservableTransfer_self}
    \leanok
    For tensors $A$ and $B$ with a common physical alphabet and possibly different
    bond dimensions, define the insertion of a $k$-site matrix $X$ by
    \begin{align}
        E_{B,A;X}(Y)=\sum_{\sigma,\tau}X_{\sigma\tau}B_\tau Y A_\sigma^*.
        \notag
    \end{align}
    Here $Y$ is a rectangular matrix from the virtual space of $A$ to that of $B$.
    For $A=B$, this is the ordinary observable insertion.
\end{definition}

\begin{lemma}[Mixed insertion identities]
    \label{thm:mixed_observable_insertion_identities}
    \lean{MPSTensor.inner_single_mixed_observable_groundSpaceMapES_single, MPSTensor.adjoint_groundSpaceMapES_observable_groundSpaceMapES_eq_mixedTransfer, MPSTensor.mixedPhysicalObservableTransfer_one, MPSTensor.mixedPhysicalObservableTransfer_appendObservable, MPSTensor.mixedPhysicalObservableTransfer_bulkObservable, MPSTensor.adjoint_groundSpaceMapES_bulkObservable_groundSpaceMapES_eq_mixedTransfer}
    \leanok
    The coefficient of $\Gamma_{A,k}^*X\Gamma_{B,k}$ between boundary units
    $(b,a)$ and $(e,c)$ is $E_{B,A;X}(|c\rangle\langle a|)_{eb}$.
    Also,
    \begin{align}
        E_{B,A;1_k}=E_{B,A}^{\,k},\qquad
        E_{B,A;X\otimes Z}=E_{B,A;X}E_{B,A;Z},
        \notag
    \end{align}
    and
    \begin{align}
        E_{B,A;1_\ell\otimes X\otimes1_r}
        =E_{B,A}^{\,\ell}E_{B,A;X}E_{B,A}^{\,r}.
        \notag
    \end{align}
\end{lemma}
\begin{proof}\leanok
    \uses{def:mixed_observable_insertion}
    Expand the boundary traces in matrix units. Splitting a physical word at a
    prescribed cut gives the composition identity, and identity insertions give
    the two ordinary transfer powers.
\end{proof}

\begin{lemma}[Decay of mixed observable sandwiches]
    \label{thm:mixed_observable_sandwich_decay}
    \lean{MPSTensor.mixedMapLM_pow_tendsto_zero_of_spectralRadius_lt_one, MPSTensor.mixedPhysicalObservableTransfer_sandwich_tendsto_zero}
    \leanok
    Suppose the rectangular mixed transfer map $E_{B,A}$ has spectral radius less
    than one. Then $E_{B,A}^{\,n}\to0$ in operator norm. For a fixed local matrix
    $X$, and along any filter with $\ell\to\infty$ and $r\to\infty$,
    \begin{align}
        E_{B,A}^{\,\ell}E_{B,A;X}E_{B,A}^{\,r}\longrightarrow0
        \notag
    \end{align}
    in operator norm.
\end{lemma}
\begin{proof}\leanok
    \uses{def:mixed_observable_insertion}
    The mixed transfer powers decay pointwise by the strict spectral-radius
    bound. Finite dimension upgrades this to operator norm convergence;
    continuity of composition then gives the inserted sandwich limit.
\end{proof}

\begin{lemma}[Decay of cross-sector boundary pullbacks]
    \label{thm:mixed_bulk_boundary_pullback_decay}
    \lean{MPSTensor.adjoint_groundSpaceMapES_bulkObservable_comp_tendsto_zero}
    \leanok
    Under the preceding strict mixed spectral-radius hypothesis, let
    $N=\ell+k+r$ and $X_{\ell,r}=1_\ell\otimes X\otimes1_r$.
    As $\ell,r\to\infty$ along any common filter,
    \begin{align}
        \bigl\|\Gamma_{A,N}^*X_{\ell,r}\Gamma_{B,N}\bigr\|\longrightarrow0.
        \notag
    \end{align}
    No injectivity of the boundary maps is required. This is a distinct-sector
    estimate, rather than an assertion of the infinite-volume spectral gap.
\end{lemma}
\begin{proof}\leanok
    \uses{thm:mixed_observable_insertion_identities,thm:mixed_observable_sandwich_decay}
    The pullback is the rectangular coefficient reshuffling of the virtual
    sandwich. This reshuffling is linear between fixed finite-dimensional
    spaces and hence continuous. Apply it to the sandwich limit.
\end{proof}

\begin{lemma}[Constructing sector states from simultaneous injectivity]
    \label{thm:normalized_multiblock_quasi_local_parent_gap}
    \lean{MPSTensor.exists_primitiveFamily_quasiLocalCommutator_gap_of_wordTupleSpanTop}
    \leanok
    Let $d>0$, and let $A_j$ have positive bond dimensions and full
    simultaneous word span at $S>0$. Let $T=\bigoplus_j\mu_jA_j$ with
    $\mu_j\neq0$, and let $h\geq0$ have kernel $G_R(T)$ at $R\geq S+1$.
    There exist normalized primitive representatives $B_j$ with
    positive-definite invariant matrices $\rho_j$ such that
    $G_N(A_j)=G_N(B_j)$ for every $j,N$. Their constructed quasi-local
    states satisfy the preceding common commutator inequality for the
    original interaction $h$, with one positive constant for all sectors,
    interval positions, and centered local observables.
    No normalized representatives or invariant matrices are given.
\end{lemma}
\begin{proof}\leanok
    \uses{thm:block_word_span_primitive_normalization,
        thm:multiblock_quasi_local_parent_commutator_gap}
    Primitive normalization preserves each block support and simultaneous
    injectivity. Hence it also preserves the combined local support and
    the parent-interaction kernel. Apply the common bound to the normalized
    family with the original positive interaction.
\end{proof}

\begin{lemma}[Quasi-local images of interior inclusions]
    \label{thm:quasi_local_bulk_inclusion}
    \lean{SpinChain.quasiLocalIntervalObservable_bulkObservable, SpinChain.quasiLocalIntervalObservable_chainWindowOperator}
    \leanok
    Let $j_{a,k}$ denote inclusion of the interval $[a,a+k)$ into the quasi-local
    algebra. Then, including empty intervals,
    \begin{align}
        j_{a-\ell,\ell+k+r}(I_\ell\otimes X\otimes I_r)=j_{a,k}(X).
        \notag
    \end{align}
    For $k>0$, the nonwrapping inclusion of $X$ at position $b$ in a chain of
    length $N$, where $b+k\leq N$, has quasi-local image $j_{a+b,k}(X)$ when the
    chain begins at $a$.
\end{lemma}
\begin{proof}\leanok
    \uses{thm:local_interval_bulk_observable,thm:quasi_local_interval_observable,lem:qca_quasi_local_observable}
    Interval coordinates identify identity extensions with finite-region
    inclusions. The quasi-local inclusions agree under enlargement of the finite
    region. The nonwrapping statement is obtained by adjoining the remaining
    sites on either side.
\end{proof}

\begin{lemma}[Fixed quasi-local support of commutator energy]
    \label{thm:quasi_local_commutator_locality}
    \lean{SpinChain.quasiLocalIntervalObservable_bulkObservable_commutator, SpinChain.quasiLocalIntervalObservable_commutator_locality}
    \leanok
    Let $R,k>0$, $t=R-1$, $w=k+2t$, and let $h$ be an arbitrary $R$-site matrix.
    Write $H_N$ for its open-chain sum and $X_N=I_\ell\otimes X\otimes I_r$, with
    $\ell,r\geq t$ and $N=\ell+k+r$. Let
    \begin{align}
        Y=(I_t\otimes X\otimes I_t)^*
        [H_w,I_t\otimes X\otimes I_t].
        \notag
    \end{align}
    Then
    \begin{align}
        j_{a-\ell,N}(X_N^*[H_N,X_N])=j_{a-t,w}(Y).
        \notag
    \end{align}
    Equivalently, putting $x=j_{a,k}(X)$ and
    $\widehat H_N=j_{a-\ell,N}(H_N)$ gives
    $x^*[\widehat H_N,x]=j_{a-t,w}(Y)$.
    Neither positivity nor a state assumption is required.
\end{lemma}
\begin{proof}\leanok
    \uses{thm:finite_open_commutator_locality,thm:quasi_local_bulk_inclusion}
    The finite-chain energy is a nonwrapping inclusion of the fixed patch
    observable $Y$. Its quasi-local image is therefore $j_{a-t,w}(Y)$.
    The equivalent form follows because the inclusions preserve products and
    adjoints.
\end{proof}

\begin{lemma}[Finite-volume commutator limit]
    \label{thm:quasi_local_commutator_energy_limit}
    \lean{SpinChain.eventually_quasiLocalIntervalObservable_bulkObservable_commutator_eq, SpinChain.tendsto_quasiLocalIntervalObservable_bulkObservable_commutator, SpinChain.tendsto_apply_quasiLocalIntervalObservable_bulkObservable_commutator, SpinChain.tendsto_apply_quasiLocalIntervalObservable_commutator}
    \leanok
    Under the preceding assumptions, if $\ell,r\to\infty$ along any common
    filter, the quasi-local energy is eventually equal to $j_{a-t,w}(Y)$ and
    hence converges to it in norm. Consequently every functional $\omega$ on the
    quasi-local algebra satisfies
    \begin{align}
        \omega(x^*[\widehat H_N,x])\longrightarrow\omega(j_{a-t,w}(Y)).
        \notag
    \end{align}
    Continuity is unnecessary here because the observable is eventually constant.
\end{lemma}
\begin{proof}\leanok
    \uses{thm:quasi_local_commutator_locality}
    Both growing margins eventually exceed $t$. Apply the fixed-support identity
    on that tail; the convergence statements follow from eventual constancy.
\end{proof}

\begin{definition}[The finite-interval ground projection]
    \label{def:mps_ground_space_projection_matrix}
    \lean{MPSTensor.groundSpaceProjectionMatrix,
        MPSTensor.isStarProjection_groundSpaceProjectionMatrix,
        MPSTensor.isStarProjection_quasiLocalIntervalObservable_groundSpaceProjectionMatrix}
    \leanok
    For an MPS tensor $A$, let $p_A(N)$ be the matrix of the orthogonal
    projection onto $G_A(N)$ in numbered coordinates on $N$ sites.
    It is a self-adjoint idempotent. Its inclusion $j_{a,N}(p_A(N))$
    into the quasi-local algebra on the interval beginning at $a$ is also
    a self-adjoint idempotent, including when $N=0$.
    \uses{def:ground_space_es, thm:quasi_local_interval_observable}
\end{definition}
\begin{proof}\leanok
    Coordinate identification and the interval inclusion preserve products
    and adjoints, hence preserve orthogonal projections.
\end{proof}

\begin{lemma}[Support compression of a positive functional]
    \label{thm:positive_functional_support_compression}
    \lean{PositiveLinearMap.apply_projection_compression_eq,
        MPSTensor.norm_quasiLocalState_interval_sub_scalar_le_projection_compression}
    \leanok
    \uses{def:qca_state_space_purity}
    Let $f$ be a positive complex-linear functional on a star-ordered
    complex algebra in which every element admits a decomposition into
    self-adjoint real and imaginary parts. If $p=p^*=p^2$ and $f(p)=f(1)$,
    then, for every $X$,
    \begin{align}
        f(pXp)=f(X).
        \notag
    \end{align}
    In particular, let $\varphi$ be a normalized positive norm-one state
    on the quasi-local algebra. For an interval projection $p$ with
    $\varphi(j_{a,N}(p))=1$, an interval observable $X$, and $c\in\mathbb C$,
    \begin{align}
        |\varphi(j_{a,N}(X))-c|\leq\|pXp-cp\|.
        \notag
    \end{align}
    No translation invariance is assumed.
\end{lemma}
\begin{proof}\leanok
    \uses{thm:quasi_local_interval_observable}
    Put $q=1-p$. Then $f(q)=f(q^*q)=0$. Cauchy--Schwarz gives
    $f(qX)=f(Xq)=0$, which proves the compression identity.
    Restrict the quasi-local state to the numbered interval matrix algebra,
    apply this identity, and use its norm-one bound on $pXp-cp$.
    The interval inclusion preserves the operator norm.
\end{proof}

\begin{lemma}[Uniqueness of a state with primitive MPS support]
    \label{thm:primitive_supported_quasi_local_state_unique}
    \lean{MPSTensor.IsPrimitiveMPS.quasiLocalState_interval_eq_insertion_of_groundSpace_support,
        MPSTensor.IsPrimitiveMPS.eq_quasiLocalExpectation_of_groundSpace_support}
    \leanok
    \uses{def:mps_ground_space_projection_matrix}
    Let $A$ be a trace-preserving primitive MPS tensor of positive physical
    and bond dimensions, with positive definite invariant matrix $\rho$. Let $\omega_A$ be its constructed quasi-local state.
    Suppose that $\varphi$ is any normalized positive norm-one continuous
    complex-linear state on the quasi-local algebra and that
    \begin{align}
        \forall\,a\in\mathbb Z,\ N\geq0:\ \varphi(j_{a,N}(p_A(N)))=1.
        \notag
    \end{align}
    Then every interval observable has the prescribed insertion expectation:
    \begin{align}
        \varphi(j_{a,k}(X))=\eta_A(X),
        \notag
    \end{align}
    and consequently $\varphi=\omega_A$. The competing state need not be
    translation invariant. This uniqueness statement is a consequence of
    primitive boundary convergence, as used in
    \cite[Sections~3 and~6]{Nachtergaele1996Spectral}; purity is a separate
    consequence of uniqueness and inheritance of support in convex
    decompositions.
\end{lemma}
\begin{proof}\leanok
    \uses{thm:positive_functional_support_compression,thm:primitive_bulk_ground_compression,thm:qca_compatible_state_extension,lem:qca_quasi_local_observable}
    Include a fixed interval observable into intervals expanded by $n$ sites
    on either side. Write $P_n$ for the ground projection of the expanded
    interval and $X_n$ for the included observable. The preceding estimate
    gives
    \begin{align}
        |\varphi(j_{a,k}(X))-\eta_A(X)|
        \leq\|P_nX_nP_n-\eta_A(X)P_n\|.
        \notag
    \end{align}
    Primitive scalar compression makes the right-hand side tend to zero.
    Thus the states agree on every interval observable. Every finite-region
    observable embeds into a containing interval, and their union is dense
    in the quasi-local algebra. Continuity proves equality on the completion.
\end{proof}

\begin{lemma}[Expectation one of MPS support projections]
    \label{thm:mps_support_projection_expectation_one}
    \lean{MPSTensor.quasiLocalExpectation_groundSpaceProjection_eq_one}
    \leanok
    Let $d>0$ and let $A$ be trace preserving. Let $\rho\ge0$ be an invariant
    virtual matrix of nonzero trace. For every $a\in\Z$ and $N\ge0$, the
    constructed state satisfies
    $\omega_A(j_{a,N}(P_{G_N(A)}))=1$.
    Primitivity and faithfulness of $\rho$ are not required.
\end{lemma}
\begin{proof}
    \uses{def:mps_ground_space_projection_matrix,thm:mps_state_local_support,thm:mps_quasi_local_expectation_state}
    \leanok
    The matrix $I-P_{G_N(A)}$ annihilates $G_N(A)$, so its translated
    expectation is zero by the local support theorem. Since $\omega_A(I)=1$,
    linearity gives expectation one on the projection.
\end{proof}

\begin{lemma}[Purity of the primitive quasi-local MPS state]
    \label{thm:primitive_quasi_local_purity}
    \lean{MPSTensor.IsPrimitiveMPS.isPureQuasiLocalState_quasiLocalExpectation}
    \leanok
    Let $d,D>0$, let $A$ be a trace-preserving primitive tensor, and let
    $\rho>0$ be its invariant virtual matrix. The completed MPS state
    $\omega_A$ is pure among all states of the quasi-local algebra.
    The states in a proposed convex decomposition need not be translation
    invariant. This proves the primitive canonical case of the purity
    assertion in \cite[Section~3]{Nachtergaele1996Spectral}; no classification
    of the multiblock ground-state face is asserted here.
\end{lemma}
\begin{proof}
    \uses{thm:qca_purity_of_supported_uniqueness,thm:mps_support_projection_expectation_one,thm:primitive_supported_quasi_local_state_unique,thm:mps_quasi_local_expectation_state,def:mps_ground_space_projection_matrix}
    \leanok
    Apply the supported-state extremality criterion to the family of
    all translated finite-interval MPS projections. They have expectation
    one in $\omega_A$, and every state satisfying the same support constraints
    equals $\omega_A$ by primitive supported-state uniqueness.
\end{proof}

\begin{lemma}[Positivity of a locally supported commutator expectation]
    \label{thm:local_commutator_expectation_positive}
    \lean{MPSTensor.observableInsertionExpectation_adjoint_commutator_nonneg}
    \leanok
    Let $A$ be a tensor, let $\rho\geq0$ be a virtual matrix, and let $H\geq0$
    act on $N$ physical sites with $G_N(A)\subseteq\ker H$. For every $N$-site
    matrix $B$,
    \begin{align}
        \eta_{A,\rho}(B^*[H,B])\geq0.
        \notag
    \end{align}
    In particular, the commutator expectation is real. Neither primitivity nor a
    spectral gap is needed.
\end{lemma}
\begin{proof}\leanok
    \uses{thm:mps_state_local_support,thm:primitive_insertion_local_state}
    The matrix $B^*BH$ annihilates $G_N(A)$, so its expectation vanishes.
    The remaining matrix $B^*HB$ is positive, and the insertion expectation is a
    positive linear functional.
\end{proof}

\begin{lemma}[Reality and positivity of sector commutator patches]
    \label{thm:sector_commutator_patch_positive}
    \lean{MPSTensor.IsParentInteraction.localCommutatorObservable_nonneg, MPSTensor.IsParentInteraction.localCommutatorObservable_im_eq_zero, MPSTensor.IsParentInteraction.quasiLocalExpectation_localCommutatorObservable_nonneg}
    \leanok
    Let $T$ be the weighted sum of tensors $A_\alpha$, with nonzero weights.
    Suppose $h\geq0$ acts on $R>0$ sites and has kernel $G_R(T)$. For a $k$-site
    matrix $X$ with $k>0$, put $t=R-1$, $w=k+2t$, and
    \begin{align}
        Y=(I_t\otimes X\otimes I_t)^*
        [H_w,I_t\otimes X\otimes I_t],
        \notag
    \end{align}
    where $H_w$ is the open-chain sum of $h$. For every block $\alpha$ and every
    virtual matrix $\rho_\alpha\geq0$,
    \begin{align}
        \eta_{A_\alpha,\rho_\alpha}(Y)\geq0,
        \qquad \operatorname{Im}\eta_{A_\alpha,\rho_\alpha}(Y)=0.
        \notag
    \end{align}
    If $A_\alpha$ is trace preserving and $\rho_\alpha$ is an invariant matrix
    of nonzero trace, the corresponding constructed quasi-local state satisfies
    $\omega_\alpha(j_{a,w}(Y))\geq0$ for every interval position $a$.
    No simultaneous injectivity bound is imposed.
\end{lemma}
\begin{proof}\leanok
    \uses{thm:local_commutator_expectation_positive,thm:open_parent_interaction_kernel_equal,thm:mps_quasi_local_expectation_state,thm:block_ground_space_le_joint,def:finite_open_interaction_matrix,lem:open_local_interaction_sum_properties,thm:open_ground_space_le_hamiltonian_kernel}
    The positive parent sum and the canonical open Hamiltonian have the same
    kernel. The canonical open Hamiltonian annihilates $G_w(T)$, which contains
    $G_w(A_\alpha)$. Apply the preceding lemma to $H_w$ and the padded observable.
    The quasi-local statement follows from exact local evaluation of the
    constructed state.
\end{proof}

\begin{lemma}[Infinite-volume commutator gap in the primitive sectors]
    \label{thm:multiblock_quasi_local_commutator_limit_gap}
    \lean{MPSTensor.exists_pos_multiblock_quasiLocalCommutator_limit_gap}
    \leanok
    Let $A=\bigoplus_{\alpha=0}^{b-1}\mu_\alpha A_\alpha$, with nonzero
    coefficients, trace-preserving primitive blocks, and faithful invariant
    density matrices. Suppose that the blocks have simultaneous injectivity
    length $S>0$, and that $h\geq0$ is a parent interaction of range
    $R\geq S+1$. Then there is $\gamma>0$ such that, for every constructed
    sector state $\omega_\alpha$ and every local observable $X$ satisfying
    $\omega_\alpha(X)=0$,
    \begin{align}
        \lim_{\ell,r\to\infty}
        \omega_\alpha\bigl(X^*[H_{[a-\ell,a+k+r)},X]\bigr)
        \geq \gamma\,\omega_\alpha(X^*X).
        \notag
    \end{align}
    Here $X$ is supported on $[a,a+k)$, with $k>0$, and the limit may be
    taken along any filter for which both margins tend to infinity.
    The limit is real and nonnegative. This proves the displayed
    infinite-volume estimate of \cite[Theorem~1.2]{Nachtergaele1996Spectral}
    for these constructed sectors at the stated sufficient range;
    identification of the entire source ground-state face is a separate result.
\end{lemma}
\begin{proof}\leanok
    \uses{thm:multiblock_quasi_local_parent_commutator_gap,thm:quasi_local_commutator_energy_limit,thm:local_commutator_expectation_positive}
    Once both margins contain $R-1$ sites, the quasi-local energy observable
    equals the inclusion of the fixed commutator patch. Thus its expectation
    is eventually constant. Local support and positivity make this expectation
    real and nonnegative. The common sector gap bounds its real part by
    $\gamma\omega_\alpha(X^*X)$, and positivity of the state makes the latter
    real as well.
\end{proof}

\begin{lemma}[Joint compression from sector column limits]
    \label{thm:sector_observable_decomposition_from_columns}
    \lean{Submodule.tendsto_norm_joint_compression_sub_sector_sum_zero}
    \leanok
    Let $J$ be finite, let $G_j(n)$ be subspaces of finite-dimensional complex
    Hilbert spaces, and let $P_j(n)$ and $P(n)$ be the projections onto
    $G_j(n)$ and $\sum_jG_j(n)$, respectively. Fix $c_j\in\C$.
    Suppose $\|\sum_jP_j(n)-P(n)\|\to0$, the operators $X_n$ are eventually
    uniformly bounded, and
    $\|P(n)X_nP_j(n)-c_jP_j(n)\|\to0$ for every $j$.
    Then $\|P(n)X_nP(n)-\sum_jc_jP_j(n)\|\to0$ along the same filter.
\end{lemma}
\begin{proof}
    \leanok
    Write $Q=\sum_jP_j$. Expanding the rightmost joint projection gives
    $PXP-\sum_jc_jP_j=PX(P-Q)+\sum_j(PXP_j-c_jP_j)$.
    Since $\|P\|\le1$,
    \begin{align}
        \label{eq:sector_joint_compression_error}
        \left\|PXP-\sum_jc_jP_j\right\|
         & \le \|Q-P\|\,\|X\|+\sum_j\|PXP_j-c_jP_j\|.
    \end{align}
    The finite sum tends to zero, as does the projection error multiplied
    by the uniform norm bound. Apply~(\ref{eq:sector_joint_compression_error}).
\end{proof}

\begin{lemma}[Sector limits of joint MPS compressions]
    \label{thm:multiblock_bulk_observable_column_limit}
    \lean{MPSTensor.bulkObservable_iSup_groundSpace_column_tendsto_scalar}
    \leanok
    Let $(A_j)_{j\in J}$ be a finite family of trace-preserving primitive
    tensors of positive bond dimensions, with faithful invariant matrices
    $\rho_j$. Suppose distinct tensors are inequivalent up to virtual gauge
    and phase. Fix a sector $\alpha$, an observable $X$ on $k>0$ sites,
    and margins $\ell_n,r_n\to\infty$ along any filter. Put
    $N_n=\ell_n+k+r_n$, let $P_j(n)=P_{G_{N_n}(A_j)}$, let $P(n)$ project
    onto $\sum_jG_{N_n}(A_j)$, and put
    $X_n=I_{\ell_n}\otimes X\otimes I_{r_n}$.
    Then $\|P(n)X_nP_\alpha(n)-\eta_{A_\alpha}(X)P_\alpha(n)\|\to0$.
\end{lemma}
\begin{proof}
    \uses{thm:multiblock_centered_full_ground_compression,thm:primitive_insertion_local_state,thm:bulk_observable_algebra}
    \leanok
    Center $X$ as $Y=X-\eta_{A_\alpha}(X)I$. Linearity and normalization
    give $\eta_{A_\alpha}(Y)=0$. Interior inclusion preserves this
    centering, and $P(n)P_\alpha(n)=P_\alpha(n)$.
    The centered full compression limit therefore gives the claimed limit.
\end{proof}

\begin{lemma}[Asymptotic sector decomposition of the compressed observable]
    \label{thm:multiblock_bulk_observable_sector_decomposition}
    \lean{MPSTensor.bulkObservable_iSup_groundSpace_sector_decomposition_tendsto_zero}
    \leanok
    Under the hypotheses and notation of
    Lemma~\ref{thm:multiblock_bulk_observable_column_limit}, with no
    distinguished sector, one has
    $\|P(n)X_nP(n)-\sum_j\eta_{A_j}(X)P_j(n)\|\to0$.
    The finite-length projections need not be mutually orthogonal.
\end{lemma}
\begin{proof}
    \uses{thm:sector_observable_decomposition_from_columns,thm:multiblock_bulk_observable_column_limit,thm:sector_projection_sum_error_decay,thm:inequivalent_block_ground_space_overlap_decay,thm:bulk_observable_algebra}
    \leanok
    The sector overlap limits and $N_n\to\infty$ imply
    $\|\sum_jP_j(n)-P(n)\|\to0$. Interior inclusion gives
    $\|X_n\|\le\|X\|$. Apply the joint-compression theorem with
    $c_j=\eta_{A_j}(X)$ and the sector column limits.
\end{proof}

\begin{lemma}[Zero energy and a support projection]
    \label{thm:positive_functional_zero_energy_support}
    \lean{PositiveLinearMap.apply_projection_eq_one_of_le_of_apply_eq_zero,
        MPSTensor.quasiLocalState_groundSpaceProjection_eq_one_of_interval_eq_zero}
    \leanok
    \uses{def:mps_ground_space_projection_matrix,def:qca_state_space_purity}
    Let $f$ be a positive normalized complex-linear functional and
    $p=p^*=p^2$. If $f(H)=0$ and $\kappa(1-p)\leq H$ for some
    $\kappa>0$, then $f(p)=1$.
    In particular, if $H\geq0$ is a finite-interval Hamiltonian with
    $\ker H=G_N(A)$, then every normalized positive quasi-local state
    satisfying $\varphi(j_{a,N}(H))=0$ also satisfies
    $\varphi(j_{a,N}(p_A(N)))=1$. No uniform gap is required.
    This is the finite-dimensional support implication used in
    \cite[Section~3]{Nachtergaele1996Spectral}.
\end{lemma}
\begin{proof}\leanok
    \uses{lem:positive-kernel-projection-comparison}
    Positivity gives $f(1-p)\geq0$, whereas monotonicity gives
    $\kappa f(1-p)\leq f(H)=0$. Thus $f(1-p)=0$ and $f(p)=1$.
    A finite positive Hamiltonian dominates a positive multiple of its
    kernel-complement projection. Restrict the quasi-local state to the
    interval matrix algebra and apply the preceding argument.
\end{proof}

\begin{lemma}[Local zero energy implies joint MPS support]
    \label{thm:quasi_local_parent_zero_energy_joint_support}
    \lean{MPSTensor.%
        quasiLocalState_groundSpaceProjection_eq_one_of_isParentInteraction_wordTupleSpanTop,
        MPSTensor.eventually_quasiLocalState_symmetric_groundSpaceProjection_eq_one}
    \leanok
    Let $A=\bigoplus_\alpha\mu_\alpha A_\alpha$ have nonzero weights and
    simultaneous injectivity length $S>0$. Let $h\geq0$ have
    $\ker h=G_R(A)$ at range $R\geq S+1$.
    If a normalized positive quasi-local state $\varphi$ satisfies
    $\varphi(j_{a,R}(h))=0$ for every $a\in\mathbb Z$, then
    \begin{align}
        \forall\,a\in\mathbb Z,\ N\geq R:\ \varphi(j_{a,N}(p_A(N)))=1.
        \notag
    \end{align}
    Consequently the same support constraint holds eventually on the
    symmetric intervals $[-n,n+1)$. The state need not be translation
    invariant, and primitive normalization of the original blocks is
    unnecessary. This is the local support consequence of the zero-energy
    condition in \cite[Theorem~1.1]{Nachtergaele1996Spectral}, at the stated
    sufficient range.
\end{lemma}
\begin{proof}\leanok
    \uses{thm:quasi_local_open_interaction_zero_expectation,thm:block_open_kernel_simultaneous_injectivity,thm:positive_functional_zero_energy_support}
    Every finite open sum has zero expectation, since each translated local
    term does. Positivity is preserved by the sum, and its kernel is exactly
    $G_N(A)$ when $N\geq R$. The preceding support lemma applies.
    Since $2n+1\geq R$ eventually, it also applies to the symmetric intervals.
\end{proof}

\begin{lemma}[Local zero energy and finite open-chain expectations]
    \label{thm:quasi_local_open_interaction_zero_expectation}
    \lean{SpinChain.quasiLocalIntervalObservable_openInteractionMatrix, SpinChain.apply_quasiLocalIntervalObservable_openInteractionMatrix_eq_zero}
    \leanok
    Let $h$ be an $R$-site matrix, where $R>0$, and let $H_N$ be its open-chain
    sum. For every interval position $a$ and every $N\geq0$,
    \begin{align}
        j_{a,N}(H_N)=\sum_{i=0}^{N-R}j_{a+i,R}(h),
        \notag
    \end{align}
    with the sum understood as empty when $N<R$.
    Consequently, any complex-linear functional $\varphi$ satisfying
    $\varphi(j_{a,R}(h))=0$ for every $a\in\mathbb Z$ also satisfies
    \begin{align}
        \varphi(j_{a,N}(H_N))=0
        \notag
    \end{align}
    for every $a$ and $N$. This applies in particular to arbitrary continuous
    states, without requiring translation invariance.
\end{lemma}
\begin{proof}\leanok
    \uses{thm:quasi_local_bulk_inclusion,def:finite_open_interaction_matrix}
    Every nonwrapping interaction window has quasi-local image $j_{a+i,R}(h)$.
    The interval inclusion preserves the finite sum. Apply linearity of
    $\varphi$ and the vanishing hypothesis to its summands.
\end{proof}

\begin{definition}[The positive finite-interval restriction]
    \label{def:positive_interval_state_restriction}
    \lean{MPSTensor.intervalStateFunctional, MPSTensor.intervalStateFunctional_apply}
    \leanok
    Let $\varphi$ be a continuous complex-linear functional on the quasi-local
    algebra that is nonnegative on every element $Y^*Y$. Its restriction to
    the numbered matrix algebra on an interval is the positive functional
    $f_{a,N}(X)=\varphi(j_{a,N}(X))$.
    \uses{thm:quasi_local_interval_observable}
\end{definition}
\begin{proof}\leanok
    Every positive matrix is a square $Y^*Y$. The interval inclusion preserves
    products and adjoints, so the restriction is positive.
\end{proof}

\begin{lemma}[Projection estimates for a finite-interval state]
    \label{thm:interval_state_projection_estimates}
    \lean{MPSTensor.intervalStateFunctional_projection_mem_unitInterval,
        MPSTensor.intervalStateFunctional_projection_eq_ofReal_re,
        MPSTensor.norm_quasiLocalState_interval_le,
        MPSTensor.quasiLocalState_interval_projection_compression_eq,
        MPSTensor.norm_quasiLocalState_interval_sub_le_projection_compression}
    \leanok
    \uses{def:positive_interval_state_restriction,def:qca_state_space_purity}
    Suppose $\varphi$ is a normalized positive norm-one quasi-local state.
    For an interval projection $p=p^*=p^2$, its expectation
    $\varphi(j_{a,N}(p))$ is real and belongs to $[0,1]$.
    Every interval observable satisfies
    $|\varphi(j_{a,N}(X))|\leq\|X\|$.
    If $\varphi(j_{a,N}(p))=1$, then
    \begin{align}
        \varphi(j_{a,N}(pXp))=\varphi(j_{a,N}(X)),\qquad
        |\varphi(j_{a,N}(X))-\varphi(j_{a,N}(Y))|
        \leq\|pXp-Y\|.
        \notag
    \end{align}
    No translation invariance is assumed.
\end{lemma}
\begin{proof}\leanok
    \uses{thm:positive_functional_support_compression,thm:quasi_local_interval_observable}
    Positivity of $p$ and $1-p$ gives the interval bounds and realness.
    The norm-one condition and the isometric interval inclusion give the
    operator-norm estimate. The support-compression identity reduces the
    final difference to the expectation of $pXp-Y$.
\end{proof}

\begin{lemma}[Compactness of finite convex weights]
    \label{thm:finite_convex_weights_compactness}
    \lean{exists_convex_weights_subsequence,
        exists_convex_weights_of_approximate_evaluations}
    \leanok
    Let $J$ be finite, and let $w_j(n)$ be real weights that eventually lie
    in $[0,1]$, with $\sum_jw_j(n)\to1$. There are a strictly increasing
    subsequence $n_s$ and weights $p_j\geq0$, $\sum_jp_j=1$, such that
    $w_j(n_s)\to p_j$ for every $j$.
    For any family of tests $x$, if
    $f(x)-\sum_jw_j(n)g_j(x)\longrightarrow0$ for every test $x$,
    then one choice of normalized weights satisfies
    $f(x)=\sum_jp_jg_j(x)$ for all tests. The test family need not be countable.
\end{lemma}
\begin{proof}\leanok
    Extract a convergent subsequence in the compact cube $[0,1]^J$.
    Continuity of the finite sum shows that the limit has total mass one.
    For each test, compare the assumed limit with the evaluation along
    this same subsequence. Uniqueness of limits gives the asserted identity.
\end{proof}

\begin{lemma}[Pure sector states and the commutator gap from the original blocks]
    \label{thm:original_blocks_pure_sector_commutator_limit_gap}
    \lean{MPSTensor.%
        exists_pure_primitiveFamily_quasiLocalCommutator_limit_gap_of_wordTupleSpanTop}
    \leanok
    Let the original block family have simultaneous injectivity length $S>0$
    and nonzero weights. For every positive parent interaction of the original
    weighted tensor at range $R\geq S+1$, there are trace-preserving primitive
    representatives, with faithful invariant matrices, preserving every
    block ground space at every length. Their completed sector states are
    pure among all quasi-local states and have a common positive
    infinite-volume commutator gap for the original interaction.
    The original blocks need not be normalized, and the invariant matrices
    and purity are conclusions. The two free margins may tend to infinity
    along any filter, as in the preceding theorem.
\end{lemma}
\begin{proof}\leanok
    \uses{thm:block_word_span_primitive_normalization,thm:multiblock_quasi_local_commutator_limit_gap,thm:primitive_quasi_local_purity}
    Normalize the blocks without changing their local ground spaces or their
    simultaneous span. Equality of the block spaces preserves the local
    kernel of the original positive interaction. Apply the common sector
    commutator theorem to the representatives and the primitive purity theorem
    to each completed sector state.
\end{proof}

\begin{lemma}[Ground-state support from the eventual kernel identity]
    \label{thm:quasi_local_zero_energy_support_eventual_kernel}
    \lean{MPSTensor.%
        eventually_quasiLocalState_symmetric_groundSpaceProjection_eq_one_of_eventual_kernel}
    \leanok
    Let $h\geq0$ have positive finite interaction range. Suppose that
    $\ker H_N(h)=G_N(A)$ for every sufficiently large $N$.
    If a normalized positive quasi-local state $\varphi$ has
    $\varphi(j_{a,R}(h))=0$ for every $a\in\mathbb Z$, then, for every sufficiently large $n$,
    \begin{align}
        \varphi\bigl(j_{-n,2n+1}(p_A(2n+1))\bigr)=1.
        \notag
    \end{align}
    This uses the eventual kernel hypothesis in
    \cite[Theorem~1.2]{Nachtergaele1996Spectral} directly. It imposes no
    local parent-kernel equality and no lower bound on the interaction
    range in terms of injectivity.
\end{lemma}
\begin{proof}\leanok
    \uses{thm:quasi_local_open_interaction_zero_expectation,thm:positive_functional_zero_energy_support}
    Every open interaction sum has zero expectation. For sufficiently large
    symmetric intervals it is positive and has the prescribed MPS kernel.
    Its finite-dimensional kernel-complement comparison therefore forces
    expectation one on the MPS support projection.
\end{proof}

\begin{definition}[The zero-energy state set of a local interaction]
    \label{def:quasi_local_parent_ground_state_face}
    \lean{MPSTensor.parentGroundStateFace}
    \leanok
    Let $d>0$ and let $h$ be an observable on $R$ consecutive sites.
    Define $\mathcal F_h$ to be the set of all states $\phi$ of the
    quasi-local algebra such that $\phi(j_{a,R}(h))=0$ for every $a\in\Z$.
    Translation invariance of $\phi$ is not required.
    For a positive parent interaction, this is its zero-energy state set,
    as in \cite[Section~1 and Theorem~1.1]{Nachtergaele1996Spectral}.
\end{definition}

\begin{lemma}[Convexity and the face property]
    \label{thm:quasi_local_parent_ground_state_face_properties}
    \lean{MPSTensor.convex_parentGroundStateFace,
        MPSTensor.isExtreme_parentGroundStateFace}
    \leanok
    The set $\mathcal F_h$ is real convex for every local observable $h$.
    If $h\ge0$, it is a face of the full state space: whenever
    $\omega=(1-t)\phi+t\psi\in\mathcal F_h$, with $0<t<1$ and
    $\phi,\psi$ states, both constituents belong to $\mathcal F_h$.
\end{lemma}
\begin{proof}
    \uses{thm:quasi_local_state_space_convex,thm:qca_state_convex_support,def:quasi_local_parent_ground_state_face}
    \leanok
    The zero-expectation equations are linear, and the state space is
    convex. For the face assertion factor $h=Y^*Y$. Each translate is
    the adjoint square $j_{a,R}(Y)^*j_{a,R}(Y)$. Vanishing of its
    expectation in a nontrivial convex combination forces vanishing in
    each constituent, by positivity.
\end{proof}

\begin{lemma}[Constructed states in the parent zero-energy face]
    \label{thm:mps_quasi_local_parent_ground_state_face_membership}
    \lean{MPSTensor.quasiLocalExpectation_mem_parentGroundStateFace_of_groundSpaceES_le_ker,
        MPSTensor.quasiLocalExpectation_mem_parentGroundStateFace_toTensorFromBlocks}
    \leanok
    Let $A$ be trace preserving and let $\rho\ge0$ be invariant with
    nonzero trace. If $G_R(A)\subseteq\ker h$, then the constructed
    state $\omega_A$ belongs to $\mathcal F_h$.
    Consequently, if $B=\bigoplus_j\mu_jA_j$, with all $\mu_j\ne0$, and
    $h$ is a positive parent interaction for $B$, every sector for which
    such trace-preserving and invariant-matrix data are given determines
    a state in $\mathcal F_h$. Neither primitivity nor faithfulness of
    the invariant matrix is required.
\end{lemma}
\begin{proof}
    \uses{thm:mps_state_local_support,thm:mps_quasi_local_expectation_state,thm:block_sum_local_ground_space,def:quasi_local_parent_ground_state_face,thm:block_ground_space_le_joint}
    \leanok
    Exact local support gives $\omega_A(j_{a,R}(h))=0$ at every position,
    while normalization and positivity make $\omega_A$ a state.
    For the weighted block tensor,
    $G_R(A_j)\subseteq\sum_iG_R(A_i)=G_R(B)=\ker h$.
    Apply the same support assertion to each sector.
\end{proof}

\begin{lemma}[Convex decomposition of a state with joint MPS support]
    \label{thm:multiblock_supported_quasi_local_state_convex_combination}
    \lean{MPSTensor.exists_convex_combination_quasiLocalExpectation_of_eventually_groundSpace_support}
    \leanok
    \uses{def:mps_ground_space_projection_matrix,def:mps_quasi_local_expectation}
    Let $(A_j)_{j=0}^{b-1}$ be finitely many trace-preserving primitive tensors
    of positive physical and bond dimensions, with positive definite invariant matrices $\rho_j$. Suppose distinct sectors are
    inequivalent up to virtual gauge and phase. Let $T$ be their weighted
    block sum with nonzero coefficients, and let $\omega_j$ be the
    constructed quasi-local state of the $j$th sector.
    Suppose $\varphi$ is a normalized positive norm-one continuous
    complex-linear state on the quasi-local algebra and that, for all
    sufficiently large $n$,
    \begin{align}
        \varphi\bigl(j_{-n,2n+1}(p_T(2n+1))\bigr)=1.
        \notag
    \end{align}
    Then there are real weights $w_j\geq0$, $\sum_jw_j=1$, such that
    \begin{align}
        \varphi=\sum_jw_j\omega_j.
        \notag
    \end{align}
    The competing state need not be translation invariant. This is the
    support-to-convex-hull implication in the argument for
    \cite[Theorem~1.1]{Nachtergaele1996Spectral}.
\end{lemma}
\begin{proof}\leanok
    \uses{thm:interval_state_projection_estimates,thm:finite_convex_weights_compactness,thm:multiblock_bulk_observable_sector_decomposition,thm:sector_projection_sum_error_decay,thm:inequivalent_block_ground_space_overlap_decay,thm:quasi_local_interval_observable,thm:mps_quasi_local_interval_expectation}
    On $I_n=[-n,n+1)$, put
    $w_j(n)=\varphi(j_{-n,2n+1}(p_{A_j}(2n+1)))$.
    Positivity gives real weights in $[0,1]$. The sum of the sector
    projections approaches the joint projection in operator norm.
    The support condition and the norm-one bound therefore give
    $\sum_jw_j(n)\to1$.

    Fix an observable $X$ on $I_m$. For $n\geq m$, include it into $I_n$
    with equal margins $n-m$. Support compression and asymptotic sector
    decomposition give
    \begin{align}
        \varphi(j_{-m,2m+1}(X))-
        \sum_jw_j(n)\eta_{A_j}(X)\longrightarrow0.
        \notag
    \end{align}
    Extract one convergent subsequence of the finite weight vectors.
    Its limit is nonnegative with total mass one. The same subsequence
    identifies all centered interval expectations with the corresponding
    convex combination. Every finite region lies in such an interval,
    and local observables are dense. Continuity proves the state identity
    on the completed algebra.
\end{proof}

\begin{lemma}[Sector zero energy from a finite open kernel]
    \label{thm:sector_local_zero_from_open_kernel}
    \lean{MPSTensor.quasiLocalExpectation_interval_eq_zero_of_open_groundSpaceES_le_ker}
    \leanok
    Let $T=\bigoplus_j\mu_jA_j$ with every $\mu_j\neq0$, and let $A_\alpha$
    be a trace-preserving block with positive invariant density
    $\rho_\alpha$ of nonzero trace, and let $\omega_\alpha$ be its constructed
    quasi-local state. Let $h$ act on $R>0$ sites. Suppose that, for some
    $N\geq R$, the joint MPS space $G_N(T)$ is contained in $\ker H_N$, where $H_N$ is
    the open-chain sum of $h$. Then $\omega_\alpha(j_{a,R}(h))=0$ for every
    interval position $a$. Positivity of $h$ is unnecessary.
\end{lemma}
\begin{proof}\leanok
    \uses{thm:quasi_local_open_interaction_zero_expectation,thm:mps_state_local_support,thm:block_ground_space_le_joint,thm:mps_quasi_local_interval_expectation}
    The sector space is contained in the joint MPS space, so
    $\omega_\alpha(j_{a,N}(H_N))=0$. Exact interval evaluation gives the
    same expectation $\eta_\alpha(h)$ for each translate of $h$. Thus
    $(N-R+1)\eta_\alpha(h)=0$. Since $N\geq R$, the coefficient is positive.
\end{proof}

\begin{lemma}[Block spaces inside the joint space]
    \label{thm:block_ground_space_le_joint}
    \lean{MPSTensor.groundSpaceES_block_le_toTensorFromBlocks}
    \leanok
    Let $T=\bigoplus_j\mu_jA_j$ with every $\mu_j\neq0$. For every block
    $\alpha$ and every length $L$, the boundary space $G_L(A_\alpha)$ is
    contained in $G_L(T)$ \cite[Lemma~commutation~(ii)]{Nachtergaele1996Spectral}.
\end{lemma}
\begin{proof}\leanok
    \uses{thm:block_sum_local_ground_space}
    With nonzero weights, the joint boundary space is the sum of the block
    boundary spaces.
\end{proof}

\begin{lemma}[Sector membership from eventual open kernels]
    \label{thm:sector_local_zero_from_eventual_open_kernel}
    \lean{MPSTensor.quasiLocalExpectation_interval_eq_zero_of_eventually_open_kernel}
    \leanok
    Under the same sector-state assumptions, suppose that
    $\ker H_N=\mathcal G_N$ for every sufficiently large $N$, where
    $\mathcal G_N$ is the joint open-boundary MPS space. Then
    $\omega_\alpha(j_{a,R}(h))=0$ for every $a$ and every sector $\alpha$.
    No injectivity condition at the interaction range is required.
    This is the sector-membership consequence of the finite-volume kernel
    hypothesis in \cite[Theorem~1.2]{Nachtergaele1996Spectral}.
\end{lemma}
\begin{proof}\leanok
    \uses{thm:sector_local_zero_from_open_kernel,def:finite_open_interaction_matrix}
    Choose a volume $N\geq R$ at which the kernel equality holds, and apply
    the finite-volume statement.
\end{proof}

\begin{lemma}[Subwindows in active-block coordinates]
    \label{lem:open_interaction_subwindow_coordinates}
    \lean{MPSTensor.cyclicCfg_join_cyclicActiveBlock,
        MPSTensor.extractWindow_join_cyclicActiveBlock,
        MPSTensor.nonwrappingSuffixStartEquiv,
        MPSTensor.periodicLocalInteractionES_apply,
        MPSTensor.periodicLocalInteractionES_conj_cyclicActiveBlock}
    \leanok
    \uses{def:cyclic_active_block_coordinates,
        def:open_local_interaction_sum}
    Suppose $W\leq N$, fix a cyclic start $s$, and let $q<W$ satisfy
    $q+R\leq W$. Replacing the range-$R$ subwindow beginning at $s+q$
    changes only the active configuration. Restriction to this subwindow
    agrees with restriction at $q$ inside the active block. When $R>0$, any local
    operator, extended by the identity, satisfies
    \[
        U_s h_{s+q}^{(N)}U_s^{-1}
        =\bigl(h_q^{(W)}\bigr)^{\oplus[d]^{N-W}}.
    \]
    Its action is obtained by applying $h$ to the restricted window with
    outside coordinates fixed. If $0<R\leq W\leq n\leq N$, translation
    $q\mapsto n-W+q$ also identifies the nonwrapping starts on $W$ sites
    with the range-$R$ starts contained in $[n-W,n)$.
\end{lemma}
\begin{proof}\leanok
    \uses{def:cyclic_active_block_coordinates}
    In active-block coordinates, subwindow replacement changes only the
    first coordinate. The start bijection has inverse subtraction of
    $n-W$. Applying the local operator separately on each complementary
    configuration gives the conjugation identity.
\end{proof}

\begin{definition}[An interaction on an embedded suffix]
    \label{def:open_suffix_interaction_sum}
    \lean{MPSTensor.openSuffixInteractionHamiltonianES}
    \leanok
    \uses{def:open_local_interaction_sum}
    For a range-$R$ operator $h$, let $H_{[n-W,n)}(h)$ be the sum of its
    nonwrapping translates in an $N$-site chain whose starts satisfy
    $n-W\leq i$ and $i+R\leq n$.
\end{definition}

\begin{lemma}[Counting positive terms in complete open windows]
    \label{lem:positive_open_window_counting}
    \lean{MPSTensor.sum_nonwrappingWindowTerms_le,
        MPSTensor.sum_openSuffixInteractionHamiltonianES_le}
    \leanok
    \uses{def:open_suffix_interaction_sum}
    Let $R\leq W$. For any positive operators $T_i$, indexed by the
    nonwrapping range-$R$ starts in an $N$-site chain,
    \[
        \sum_{n=W}^{N}\ \sum_{\substack{n-W\leq i\\i+R\leq n}}T_i
        \leq (W-R+1)\sum_i T_i.
    \]
    In particular, for a positive interaction $h$,
    $\sum_{n=W}^{N}H_{[n-W,n)}(h)
        \leq(W-R+1)H_N^{\mathrm{open}}(h)$.
    This is the counting argument of condition C1 in
    \cite{Nachtergaele1996Spectral}, applied to arbitrary positive terms.
\end{lemma}
\begin{proof}\leanok
    \uses{def:open_suffix_interaction_sum}
    A term beginning at $i$ can occur only for
    $i+R\leq n\leq i+W$, an interval containing $W-R+1$ integers.
    Reverse the order of summation and use positivity of each term.
\end{proof}

\begin{lemma}[Transfer of a finite-window comparison]
    \label{lem:open_interaction_suffix_comparison}
    \lean{MPSTensor.openSuffixInteractionHamiltonianES_conj_cyclicActiveBlock,
        MPSTensor.openSuffixInteractionHamiltonianES_comparison_of_local}
    \leanok
    \uses{lem:open_interaction_subwindow_coordinates,
        def:open_suffix_interaction_sum}
    Suppose $0<R\leq W\leq n\leq N$. The suffix Hamiltonian is unitarily
    equivalent to $H_W^{\mathrm{open}}(h)$ acting independently on the
    complementary configurations. If
    $\kappa p_W(A)\leq H_W^{\mathrm{open}}(h)$, then the same inequality
    holds after embedding both sides on the suffix $[n-W,n)$.
    Here $p_W(A)$ is the canonical complementary ground-space projection;
    no sign condition on $\kappa$ or positivity condition on $h$ is needed
    for this transfer of order.
\end{lemma}
\begin{proof}\leanok
    \uses{lem:open_interaction_subwindow_coordinates}
    Sum the conjugation identity over the corresponding start bijection.
    Extension by the identity and conjugation by an isometry preserve order.
\end{proof}

\begin{lemma}[Open-chain comparison for an arbitrary positive interaction]
    \label{thm:positive_open_interaction_range_comparison}
    \lean{MPSTensor.openInteractionHamiltonianES_comparison_of_local,
        MPSTensor.openInteractionHamiltonianES_gap_of_long_gap_at_length}
    \leanok
    \uses{lem:positive_open_window_counting,
        lem:open_interaction_suffix_comparison,
        lem:open_parent_range_comparison}
    Let $0<R\leq W\leq N$, let $h$ be positive, and suppose
    $\kappa p_W(A)\leq H_W^{\mathrm{open}}(h)$. Then
    \[
        \kappa H_{N,W}^{\mathrm{open}}(A)
        \leq (W-R+1)H_N^{\mathrm{open}}(h).
    \]
    If additionally $\kappa,\gamma>0$, both global kernels equal $G_N(A)$,
    and $\gamma\lVert v\rVert\leq
        \lVert H_{N,W}^{\mathrm{open}}(A)v\rVert$ for $v\perp G_N(A)$, then
    \[
        \frac{\kappa\gamma}{W-R+1}\lVert v\rVert
        \leq\lVert H_N^{\mathrm{open}}(h)v\rVert
        \qquad(v\perp G_N(A)).
    \]
    No local parent-kernel equality is assumed for $h$. This is a
    finite-window consequence of condition C1 and the comparison argument
    in Section~6 of \cite{Nachtergaele1996Spectral}.
\end{lemma}
\begin{proof}\leanok
    \uses{lem:positive_open_window_counting,
        lem:open_interaction_suffix_comparison}
    Sum the embedded local comparisons. Each full range-$W$ parent term
    occurs once, whereas a range-$R$ term occurs at most $W-R+1$ times.
    For the gap, compare the positive operators
    $\kappa H_{N,W}^{\mathrm{open}}(A)$ and
    $(W-R+1)H_N^{\mathrm{open}}(h)$, whose kernels agree, and divide by
    the positive window multiplicity.
\end{proof}

\begin{lemma}[Uniform finite-chain gap from an eventual kernel identity]
    \label{thm:eventual_kernel_open_interaction_gap}
    \lean{MPSTensor.exists_openInteractionHamiltonianES_gap_of_eventual_kernel,
        MPSTensor.exists_openInteractionMatrix_gap_of_eventual_kernel}
    \leanok
    Let $A_j$ be a finite family of normalized primitive MPS tensors with
    faithful fixed density matrices, pairwise inequivalent under gauge and
    phase, and let $B=\bigoplus_j\mu_jA_j$, where every $\mu_j\ne0$.
    Let $h\geq0$ be an interaction of range $R\geq1$.
    Suppose that $\ker H^{\mathrm{open}}_{h,N}=\mathcal G_N(B)$
    for all sufficiently large $N$. Then there exists $\gamma>0$ such that
    \begin{align}
        \gamma\|v\|\leq\|H^{\mathrm{open}}_{h,N}v\|
        \notag
    \end{align}
    for every $N\in\mathbb N$ and every
    $v\perp\ker H^{\mathrm{open}}_{h,N}$.
    No condition on the local kernel of $h$, or lower bound on $R$
    beyond $R\geq1$, is required. This is the finite-chain conclusion of
    \cite[Theorem~1.2]{Nachtergaele1996Spectral}, under its eventual
    finite-interval kernel hypothesis, for the primitive-sector representation;
    the passage from more general periodic presentations is separate
    \cite{gap:nachtergaele96_infinite_volume_ground_projection}.
\end{lemma}
\begin{proof}\leanok
    \uses{thm:primitive_block_sum_open_gap_all_lengths,
        thm:ph_primitive_blocks_open_kernel,
        lem:parent-finite-interval-comparison,
        thm:positive_open_interaction_range_comparison}
    Choose a canonical interaction range $W=2p$ above the original range
    and both kernel thresholds, with canonical gap $\delta>0$ uniform in
    the volume. Finite dimensionality and positivity give
    $\kappa P_{\mathcal G_W(B)^\perp}\leq H^{\mathrm{open}}_{h,W}$
    for some $\kappa>0$. For $N\geq W$, the finite-window comparison gives
    \begin{align}
        \kappa H^{\mathrm{open}}_{N,W}(B)
        \leq (W-R+1)H^{\mathrm{open}}_{h,N},
        \notag
    \end{align}
    The kernels coincide for these volumes, so the transferred gap is
    $\kappa\delta/(W-R+1)$. Taking a minimum with the positive gaps
    on the kernel complements at the finitely many volumes $N<W$
    gives one constant for every volume.
\end{proof}

\begin{theorem}[Classification of the zero-energy state face for primitive sectors]
    \label{thm:primitive_multiblock_quasi_local_ground_state_classification}
    \lean{MPSTensor.mem_parentGroundStateFace_iff_convex_combination_of_eventually_open_kernel}
    \leanok
    \uses{def:quasi_local_parent_ground_state_face}
    Let $T=\bigoplus_{j=0}^{b-1}\mu_jA_j$, with $\mu_j\ne0$, be a finite
    weighted sum of inequivalent trace-preserving primitive tensors with
    faithful invariant density matrices. Write $\omega_j$ for their
    constructed quasi-local states. Let $h\geq0$ have range $R\geq1$,
    and suppose $\ker H_N^{\mathrm{open}}(h)=G_N(T)$ for all sufficiently
    large $N$. Then, for every normalized positive state $\varphi$,
    \begin{align}
        \varphi\in\mathcal F_h
        \quad\Longleftrightarrow\quad
        \exists (w_j)_j:\ w_j\geq0,\quad\sum_jw_j=1,\quad
        \varphi=\sum_j w_j\omega_j.
        \notag
    \end{align}
    No translation invariance of $\varphi$ is assumed. This is the
    ground-state classification associated with
    \cite[Theorems~1.1--1.2]{Nachtergaele1996Spectral} for the stated
    primitive sector family. More general periodic presentations remain
    separate \cite{gap:nachtergaele96_infinite_volume_ground_projection}.
\end{theorem}
\begin{proof}\leanok
    \uses{thm:quasi_local_zero_energy_support_eventual_kernel,thm:multiblock_supported_quasi_local_state_convex_combination,thm:quasi_local_parent_ground_state_face_properties,thm:sector_local_zero_from_eventual_open_kernel}
    Zero local energy and the eventual kernel identity give expectation
    one on the joint support projections of sufficiently large intervals.
    The supported-state decomposition gives the convex coefficients.
    Conversely, each sector state has zero local interaction expectation:
    its expectation of a sufficiently long open sum is zero, and every
    summand has the same expectation. Convexity preserves this property.
\end{proof}

\begin{corollary}[Pure ground states are precisely the primitive sectors]
    \label{thm:primitive_multiblock_pure_ground_state_classification}
    \lean{MPSTensor.isPure_mem_parentGroundStateFace_iff_sector_of_eventually_open_kernel}
    \leanok
    Under the preceding hypotheses, a state $\varphi$ is a pure element of
    $\mathcal F_h$ if and only if $\varphi=\omega_j$ for some $j$.
    Purity is extremality among all states of the quasi-local algebra.
\end{corollary}
\begin{proof}\leanok
    \uses{thm:primitive_multiblock_quasi_local_ground_state_classification,thm:finite_pure_state_decomposition,thm:primitive_quasi_local_purity}
    A pure state in a finite convex combination equals one of the
    constituent states. Each constructed primitive sector state is pure
    and has zero local energy.
\end{proof}

\begin{lemma}[Positive commutator expectation from eventual joint support]
    \label{thm:eventual_kernel_sector_commutator_positive}
    \lean{MPSTensor.quasiLocalExpectation_localCommutatorObservable_nonneg_of_eventual_kernel}
    \leanok
    Let $h\geq0$ have range $R>0$, and suppose that the kernels of its
    sufficiently long open-chain Hamiltonians equal the joint boundary spaces
    of a weighted tensor family with nonzero weights. Let $\omega_\alpha$ be
    a constructed sector state, with a trace-preserving tensor and a positive
    invariant matrix of nonzero trace. For $X$ on $k>0$ sites,
    \begin{align}
        \omega_\alpha\bigl(j_{a-(R-1),\,2(R-1)+k}
        (X^*[H_{\mathrm{patch}},X])\bigr)\geq0.
        \notag
    \end{align}
    Neither primitivity nor faithfulness of the invariant matrix is needed
    for this assertion. The local commutator patch has real expectation.
\end{lemma}
\begin{proof}\leanok
    \uses{thm:local_commutator_expectation_positive,thm:quasi_local_commutator_locality,thm:mps_quasi_local_interval_expectation}
    Choose one sufficiently large volume containing both interaction
    margins. Its Hamiltonian annihilates the sector boundary space and is
    positive. The expectation of $B^*[H,B]$ therefore equals that of
    $B^*HB$ and is nonnegative. The quasi-local commutator identity identifies
    this expectation with the fixed local patch.
\end{proof}

\begin{lemma}[One finite-chain and local commutator gap for all primitive sectors]
    \label{thm:eventual_kernel_sector_local_gap}
    \lean{MPSTensor.exists_pos_multiblock_localCommutator_gap_of_eventual_kernel,
        MPSTensor.exists_pos_multiblock_quasiLocalCommutator_gap_of_eventual_kernel}
    \leanok
    Let $A_1,\ldots,A_b$ be pairwise gauge-phase inequivalent normalized
    primitive tensors with faithful invariant matrices, and let all weights
    be nonzero. Let $h\geq0$ have arbitrary range $R>0$. Suppose that
    $\ker H_N^{\mathrm{open}}(h)$ equals the joint boundary-condition space
    for all sufficiently large $N$. Then there is $\gamma>0$ such that, for every $N$ and every
    $v\perp\ker H_N^{\mathrm{open}}(h)$,
    \begin{align}
        \gamma\lVert v\rVert\leq\lVert H_N^{\mathrm{open}}(h)v\rVert,
        \notag
    \end{align}
    and, for every sector, position, and centered observable $X$ on $k>0$
    sites,
    \begin{align}
        \gamma\omega_\alpha(X^*X)
        \leq\omega_\alpha(X^*[H_{\mathrm{patch}},X]).
        \notag
    \end{align}
    The same inequality holds for the finite-interval insertion expectations.
    No local parent-kernel equality or lower bound on $R$ in terms of an
    injectivity length is assumed. This is the primitive-sector specialization
    of the argument in \cite[Theorem~1.2 and Section~6]{Nachtergaele1996Spectral};
    the more general GVBS presentations are recorded in
    \cite{gap:nachtergaele96_infinite_volume_ground_projection}.
\end{lemma}
\begin{proof}\leanok
    \uses{thm:eventual_kernel_open_interaction_gap,thm:multiblock_local_open_interaction_commutator_gap,thm:mps_quasi_local_interval_expectation}
    Choose the uniform finite-chain constant. Pull the eventual kernel
    equality along growing buffered volumes. Apply the full joint-ground
    projection decay and the finite-volume commutator estimate in each sector.
    Evaluation on consecutive intervals is the insertion expectation.
\end{proof}

\begin{lemma}[A common finite-chain and infinite-volume sector gap]
    \label{thm:eventual_kernel_sector_literal_gap}
    \lean{MPSTensor.exists_pos_multiblock_quasiLocalCommutator_limit_gap_of_eventual_kernel}
    \leanok
    Under the primitive-family hypotheses of
    Lemma~\ref{thm:eventual_kernel_sector_local_gap}, choose the same
    $\gamma>0$ for the finite-chain bound. Let $\ell,r$ tend to infinity
    along any filter. For every constructed sector state and centered local
    $X$ on $k>0$ sites, the finite-volume energies
    \begin{align}
        \omega_\alpha\bigl(X^*[H_{[a-\ell,\,a+k+r)}(h),X]\bigr)
        \notag
    \end{align}
    converge to a real number $e$ satisfying
    $\gamma\omega_\alpha(X^*X)\leq e$. Thus one constant gives both the
    all-volume finite-chain norm gap and the infinite-volume sector bound.
    This states the primitive-sector case explicitly; classification of all
    pure zero-energy states is a separate assertion.
\end{lemma}
\begin{proof}\leanok
    \uses{thm:eventual_kernel_sector_local_gap,thm:eventual_kernel_sector_commutator_positive,thm:quasi_local_commutator_locality}
    Once both margins exceed $R-1$, locality makes the commutator observable
    equal to its fixed patch. Its expectation is real and nonnegative by
    eventual support. Combine this finite-volume limit with the local bound using
    the same finite-chain constant.
\end{proof}

\begin{lemma}[Ground-state classification after primitive normalization]
    \label{thm:normalized_quasi_local_ground_state_classification}
    \lean{MPSTensor.%
        exists_normalized_parentGroundStateFace_classification_of_eventually_open_kernel}
    \leanok
    Let $A_j$ be a finite family with positive bond dimensions whose
    word tuples at some length $S>0$ span the full product matrix algebra.
    This is the simultaneous spanning condition of
    \cite[Section~IV.C]{Cirac2021Matrix}.
    Let $T_A=\bigoplus_j\mu_jA_j$, with every $\mu_j\ne0$.
    Suppose that $h\geq0$ has range $R\geq1$ and that
    $\ker H^{\mathrm{open}}_{h,N}=G_N(T_A)$ for all sufficiently large $N$.
    There exist primitive trace-preserving tensors $B_j$ with faithful
    invariant matrices $\rho_j$ such that their word tuples still span the
    full product algebra at length $S$. Write $T_B=\bigoplus_j\mu_jB_j$. Then
    \begin{align}
        G_N(A_j) & =G_N(B_j), &
        G_N(T_A) & =G_N(T_B).
        \notag
    \end{align}
    These equalities hold for every block $j$ and every length $N$.
    Let $\omega_j$ be the quasi-local state constructed from $B_j,\rho_j$.
    For every continuous complex-linear functional $\varphi$ on the
    quasi-local algebra,
    \begin{align}
        \varphi\in\mathcal F_h
         & \quad\Longleftrightarrow\quad
        \exists (w_j)_j:\ w_j\geq0,\quad\sum_jw_j=1,\quad
        \varphi=\sum_jw_j\omega_j,\notag \\
        \varphi\in\mathcal F_h\ \text{and }\varphi\text{ is pure}
         & \quad\Longleftrightarrow\quad
        \exists j:\ \varphi=\omega_j.\notag
    \end{align}
    The original interaction range need not exceed $S$.
    This combines the normalization in
    \cite[equations~(3.1)--(3.2b)]{Nachtergaele1996Spectral}
    with the ground-state classification of
    \cite[Theorems~1.1--1.2]{Nachtergaele1996Spectral}, for the original
    simultaneously spanning family. The passage from more general periodic
    presentations is separate
    \cite{gap:nachtergaele96_infinite_volume_ground_projection}.
\end{lemma}
\begin{proof}\leanok
    \uses{thm:block_word_span_primitive_normalization,
        thm:block-word-span-separation,thm:ph_block_replacement_parent_equality,
        thm:primitive_multiblock_quasi_local_ground_state_classification,
        thm:primitive_multiblock_pure_ground_state_classification}
    Normalize each block by a nonzero rescaling and a virtual gauge.
    These operations preserve every block ground space and the simultaneous
    word span. The weighted joint ground spaces therefore coincide at every
    length, and the eventual kernel identity transfers to $T_B$.
    Simultaneous spanning excludes gauge-phase equivalence between distinct
    normalized blocks. Apply the convex-combination and pure-sector
    classifications to this normalized family and the unchanged interaction.
\end{proof}

\begin{theorem}[Finite and infinite-volume gaps in every pure ground state]
    \label{thm:eventual_kernel_pure_ground_state_gap}
    \lean{MPSTensor.exists_pos_pure_ground_state_commutator_limit_gap_of_eventual_kernel}
    \leanok
    Let $T=\bigoplus_j\mu_jA_j$ be a finite weighted family of inequivalent
    trace-preserving primitive tensors with faithful invariant density matrices
    and nonzero coefficients. Let $h\geq0$ have range $R\geq1$, and suppose
    $\ker H_N^{\mathrm{open}}(h)=G_N(T)$ for every sufficiently large $N$.
    Let both free margins tend to infinity. Then there is one $\gamma>0$
    such that
    \begin{align}
        \gamma\|v\|\leq\|H_N^{\mathrm{open}}(h)v\|
        \qquad\bigl(N\geq0,\ v\perp\ker H_N^{\mathrm{open}}(h)\bigr),
        \notag
    \end{align}
    and, for every pure $\eta\in\mathcal F_h$ and every local observable
    $X$ with $\eta(X)=0$, the finite-volume limit exists and satisfies
    \begin{align}
        \lim_{M\to-\infty,\,N\to+\infty}
        \eta\bigl(X^*[H_{[M,N]}(h),X]\bigr)
        \geq\gamma\eta(X^*X).
        \notag
    \end{align}
    The margins may diverge along an arbitrary filter. This is the pair
    of conclusions in \cite[Theorem~1.2]{Nachtergaele1996Spectral} for
    the stated primitive family; the original range has no injectivity
    lower bound. General periodic presentations remain separate
    \cite{gap:nachtergaele96_infinite_volume_ground_projection}.
\end{theorem}
\begin{proof}\leanok
    \uses{thm:eventual_kernel_sector_literal_gap,thm:primitive_multiblock_pure_ground_state_classification}
    The sector gap theorem gives one constant for both conclusions.
    The pure-state classification identifies each pure zero-energy state
    with a constructed sector, to which the same finite-volume limit and bound apply.
\end{proof}

\begin{lemma}[The pure ground-state gap from an original spanning family]
    \label{thm:original_family_eventual_kernel_pure_ground_state_gap}
    \lean{MPSTensor.exists_pos_pure_ground_state_commutator_limit_gap_of_wordTupleSpanTop}
    \leanok
    The preceding finite and infinite-volume conclusions also hold when
    the original blocks are only assumed to have full simultaneous word
    span at some positive length $S$. Their normalization and invariant
    density matrices are derived internally. No condition $R\geq S+1$
    is required; positivity and the eventual open-chain kernel identity
    remain the hypotheses on the original interaction.
\end{lemma}
\begin{proof}\leanok
    \uses{thm:normalized_quasi_local_ground_state_classification,thm:eventual_kernel_pure_ground_state_gap}
    Normalize the blocks, preserving every block ground space and every
    weighted joint ground space. The eventual kernel condition transfers
    to this family. Its simultaneous span excludes gauge-and-phase
    equivalence of different blocks. Apply the preceding theorem.
\end{proof}

\begin{lemma}[Existence of a positive interaction for a primitive family]
    \label{thm:primitive_ground_state_interaction_exists}
    \lean{MPSTensor.exists_positive_parent_interaction_of_isPrimitiveMPS}
    \leanok
    \uses{def:open_local_interaction_sum}
    Let $T=\bigoplus_{j=0}^{b-1}\mu_jA_j$, with nonzero weights, be a finite
    family of pairwise gauge-phase inequivalent trace-preserving primitive
    tensors of positive bond dimensions. Suppose each tensor has a faithful
    invariant matrix. There exist $R>0$ and a positive interaction $h$ on
    $R$ sites such that, for every $N\geq R$,
    \begin{align}
        \ker H_N^{\mathrm{open}}(h)=G_N(T).
        \notag
    \end{align}
    Neither $R$ nor a simultaneous injectivity length is given.
    The inequality $N\geq R$ is stronger than the inclusive-interval bound
    $N\geq R+1$ in
    \cite[Theorem~1.1]{Nachtergaele1996Spectral}.
\end{lemma}
\begin{proof}\leanok
    \uses{thm:ph_primitive_blocks_open_kernel,rem:euclidean_local_projector_ingredients}
    Choose the positive open-kernel threshold $R$ and let $h$ be the
    orthogonal projection onto $G_R(T)^\perp$. Its open sum is the canonical
    range-$R$ parent Hamiltonian. The threshold theorem gives the stated
    kernel equality, including the single-window volume $N=R$.
\end{proof}

\begin{lemma}[An interaction realizing the entire primitive-sector ground-state face]
    \label{thm:primitive_ground_state_face_interaction_exists}
    \lean{MPSTensor.exists_positive_parent_interaction_groundStateFace_of_isPrimitiveMPS}
    \leanok
    For the primitive family of the preceding theorem, one can choose a
    single $R>0$ and $h\geq0$ such that the open kernels equal $G_N(T)$
    for every $N\geq R$, and hence for all sufficiently large $N$, and
    \begin{align}
        \mathcal F_h
        =\left\{\sum_{j=0}^{b-1} w_j\omega_j:
        w_j\geq0,\ \sum_jw_j=1\right\}.
        \notag
    \end{align}
    Here $\omega_j$ is the constructed quasi-local state of $A_j$.
    The pure states in $\mathcal F_h$ are precisely the states $\omega_j$.
    Purity is extremality among all normalized positive quasi-local states;
    no translation invariance of a state being classified is assumed.
    This is the stated primitive-family case of the existence assertion in
    \cite[Theorem~1.1 and its finite-family consequence]{Nachtergaele1996Spectral}.
    More general periodic GVBS presentations remain separate
    \cite{gap:nachtergaele96_infinite_volume_ground_projection}.
\end{lemma}
\begin{proof}\leanok
    \uses{thm:primitive_ground_state_interaction_exists,thm:primitive_multiblock_quasi_local_ground_state_classification,thm:primitive_multiblock_pure_ground_state_classification}
    The all-volume kernel bound implies the eventual kernel hypothesis.
    Apply the two classification theorems to the same positive canonical
    interaction. They characterize the entire zero-energy face and its pure
    states without an invariance assumption on the state.
\end{proof}
